\documentclass[11pt]{article}
\usepackage[T1]{fontenc}
\usepackage[utf8]{inputenc}
\usepackage{lmodern,amsmath,amssymb,amsthm,mathtools,mathrsfs}
\usepackage[a4paper,margin=28mm]{geometry}
\usepackage{microtype,booktabs,array,longtable,enumitem}
\usepackage[hidelinks]{hyperref}
\usepackage{xurl}
\setlength{\emergencystretch}{3em}
\setlist[enumerate]{leftmargin=*,itemsep=.3em}
\newtheorem{theorem}{Theorem}[section]
\newtheorem{proposition}[theorem]{Proposition}
\newtheorem{lemma}[theorem]{Lemma}
\newtheorem{corollary}[theorem]{Corollary}
\theoremstyle{definition}
\newtheorem{definition}[theorem]{Definition}
\newtheorem{example}[theorem]{Example}
\theoremstyle{remark}
\newtheorem{remark}[theorem]{Remark}
\newcommand{\Q}{\mathbf Q}
\newcommand{\Z}{\mathbf Z}
\newcommand{\C}{\mathbf C}
\newcommand{\A}{\mathcal A}
\newcommand{\cl}{\operatorname{cl}}
\newcommand{\Spec}{\operatorname{Spec}}
\newcommand{\Gal}{\operatorname{Gal}}
\newcommand{\MT}{\operatorname{MT}}
\newcommand{\Nm}{\operatorname{N}}
\newcommand{\tdeg}{\operatorname{tdeg}}
\newcommand{\workbenchurl}{https://github.com/KokunoYumeto/algebraic-cycles-and-motives-workbench}
\newcommand{\recordurl}{https://github.com/KokunoYumeto/mathematics-workbenches}
\title{An effective one-prime criterion for Weil classes\\over $\Q(i)$ and $\Q(\omega)$}
\author{}
\date{}
\begin{document}
\maketitle
\begin{abstract}
Br\"unjes--Serp\'e's arithmetic lifting criterion and the bounded-cycle
criterion for the Hodge conjecture reduce algebraicity to a uniform
degree bound. We prove a one-prime implication with a verifiable
arithmetic certificate and a polynomial discriminant threshold for
relative-Weyl witnesses, with explicit unitary datum constants. An integer formula bounds
the degree of the detector locus in terms of the cycle degree.
The quantitative inputs are Jaswal's unrefereed effective
Andr\'e--Oort theorem and the zero-free half-planes of
Cureton--Thorner or Hu. Qualitative Andr\'e--Oort for these Hodge-type families follows
from Tsimerman's published theorem for $\mathcal A_g$. Alp\"oge's asserted Dirichlet
GRH gives a stronger half-plane and explicit splitting primes for
abelian fields; the full relative-Weyl witnesses here have
nonabelian splitting fields. Milne's finite-field Tate theorem,
conditional here on the OpenAI preprints of 6 October 2026 and their
unrefereed human write-ups, supplies algebraicity but no degree bound.
CM height estimates sharpen the point threshold, with splitting-field
exponents $768$ and $10\,569\,646\,080$ in the total degree of the
base locus, up to logarithmic factors. We also give an ordinary-reduction
discriminant restriction and arithmetic degree-escape certificates.
The Gaussian eightfold detector remains unconstructed; no Hodge
conjecture for that family is proved. The mathematics and text were
produced by GPT-6 Astra, with separate same-model checks described
below; no human has checked this work.
\end{abstract}

\section{The Hodge conjecture and a test at one prime}

Weil asked whether the exceptional Hodge classes furnished by an
imaginary quadratic action on an abelian variety are algebraic.
For a polarized abelian variety $A$ of dimension $2n$ with an action
of $K=\Q(i)$ or $\Q(\omega)$ of signature $(n,n)$, the two-dimensional
rational space in question is
\[
 W_K(A)=\bigwedge_K^{2n}H^1(A,\Q)\ \subset\ H^{2n}(A,\Q).
\]
Here $\omega^2+\omega+1=0$, and the polarization is compatible with
complex conjugation on $K$. A \emph{detector} is an irreducible
$n$-dimensional subvariety whose cycle class has nonzero projection
to $W_K(A)$. Its degree is $\int_Z\theta^n$, without division by
$n!$. Write $\delta_W(A,\theta)$ for the least such degree, or
$+\infty$ if there is no detector. One detector makes the whole
Weil plane algebraic; we prove the elementary extraction needed
for this assertion below.

The bounded-cycle criterion in \cite[Family criterion and section on
one finite-field detector]{Bounded} starts from proper parameter
spaces for cycles. It says, in particular, that a uniform bound on
detector degrees on CM fibres of a fixed unitary family suffices on
every complex fibre. Its arithmetic form requires one bound for
infinitely many good reductions of every CM fibre. Br\"unjes and
Serp\'e prove the underlying uniform-complexity lifting principle
in \cite[Theorem 3.4, p.~8]{BrunjesSerpe}.

There are two distinct ways in which one prime can suffice. A
cycle outside an explicitly certified exceptional set lifts to its
\emph{fixed} characteristic-zero fibre. A cycle at a point where
the cycle parameter space is flat over the \emph{moduli base}
forces a detector throughout that base. The latter flatness is a
geometric condition and does not follow from ordinarity.

\paragraph{Arithmetic inputs and their scope.}
The closed-image, flatness, cycle-class and CM-specialization inputs
are established algebraic geometry. For the Hodge-type families here,
qualitative Andr\'e--Oort follows from Tsimerman's theorem for
$\mathcal A_g$ \cite[Theorem~1.3, p.~380]{Tsimerman}, by the
finite Shimura map to Siegel moduli. Pila--Shankar--Tsimerman, with
Esnault--Groechenig, prove Andr\'e--Oort in full generality
\cite[Theorem~1.1]{PST}, in a preprint. Jaswal's effective refinement
\cite[Theorem 1.1(ii)]{Jaswal} is a separate unrefereed input, whose
Hexagon record declares primarily AI-generated text and no human
understanding of its mathematics. Cureton--Thorner's
\cite[Theorem 1.1]{CuretonThorner} supplies the parameter
$\delta=1/32$ in its zero-free hypothesis. Hu's
\cite[Theorem 1.2]{Hu} supplies $\delta=1/8$. These recent analytic
inputs are named whenever they are used. Alp\"oge's announced
Dirichlet GRH \cite[p.~1]{Alpoge} supplies $\delta=1/2$ for abelian
splitting fields, and, with Bach--Sorenson, a prime bound in
arithmetic progressions. We use the asserted Dirichlet GRH as an
input; this article does not verify its formal development.

Milne's report \cite[Theorem 3.1]{MilneReport}, conditional here on
the two OpenAI preprints of 6 October 2026 and their unrefereed
human write-ups, gives the full Tate conjecture for abelian varieties
over finite fields. Under those inputs a specialized Weil class is
algebraic. That conclusion contains no bound on the degree of a
representing cycle. Thus the arithmetic bounded-cycle criterion
asks for \emph{uniformity in the Tate theorem}. None of our
one-prime implications uses the Tate theorem to manufacture a
degree bound.

\subsection{The main implications}

Fix a geometrically connected fine unitary Shimura variety $S/F$,
with neat level of product form, and a polarized abelian scheme
with labelled integral $K$-action. Its adjoint group is assumed
$\Q$-simple, with absolute root systems of type $A$. Degrees on
the Baily--Borel compactification use the fixed polarization in
the effective Andr\'e--Oort input. A special point is called
\emph{relative Weyl} when its adjoint Mumford--Tate torus is
maximal and the Galois action on its character space contains the
Weyl group of this unitary datum. Let
\[
 \mathcal D_D=\{s\in S:\delta_W(s)\leq D\}.
\]
For its geometric components $Z_j$, choose an integer
\begin{equation}
 H_D\ \geq\ \sum_j\max\{1,\lceil\deg_{\rm BB}\overline Z_j\rceil\}.
 \label{eq:total-degree}
\end{equation}
The closure is taken in the Baily--Borel compactification. This
convention controls both degrees and the number of components.
Let $c_S(\delta,h,d)$ be the constant of Jaswal's theorem for the
fixed Shimura datum, enlarged to be nondecreasing in the integer
arguments $h,d$ by taking their finite running maximum.

\begin{theorem}\label{thm:explicit-theta-chow-bound}
Let $g=2n\in\{4,8\}$, put $s=n^2$, and let $S$ be a geometrically
connected component of the principal-level-$3$ unitary family for
$K=\mathbf Q(i)$ or $\mathbf Q(\sqrt{-3})$, on the prescribed
principally polarized integral symplectic lattice.  The level is taken
in the underlying representation of rank $2g$.  For $D\geq1$ an
integer, put
\begin{align*}
 N_\theta&=8^g-1,& k&=8^nD,& c&=\binom{N_\theta+k}{k}^{n+1},\\
 \ell&=(N_\theta+1)^k,& h&=(n+1)k,& t&=\ell(h-1)+1,\\
 b&=\binom{\ell+t-1}{t},& e&=2b,& R&=(N_\theta+1)c-1,
\end{align*}
and
\[
 I_g=3^{g^2}\prod_{j=1}^{g}(3^{2j}-1),\qquad
 H_g(D)=I_g\,g^s\,k\,e^R.
\]
Then $H_D=H_g(D)$ satisfies \eqref{eq:total-degree} for the
reduced detector locus $\mathcal D_D\subset S$, including the
rounded positive degree assigned to each geometric component.  The same formula
works for each geometrically connected component, and for either of
the two fields $K$.
\end{theorem}

This supplies the degree estimate needed in the forcing criterion
without an unevaluated projective-model constant.  The proof uses
Mumford's theta equations and an explicit elimination argument for
relative Chow forms; it does not construct a detecting cycle.
For the known degree-$48$ even Gaussian and degree-$36$ Eisenstein
fourfold bounds, the detector locus is the entire base, and the
smaller choice
\[
 H_D=131569513308979200\cdot4^{4095}
\]
is available by Corollary~\ref{cor:theta-degree-whole-base}.

The main numerical conventions are as follows.  The letters used
for local projective calculations are specified again in those
calculations.
\begin{center}
\begin{tabular}{ll}
$m=g=2n$ & dimension of an abelian fibre, here $4$ or $8$\\
$s=n^2$ & dimension of the unitary base\\
$D$ & degree of a cycle in a fibre\\
$H_D$ & rounded total degree of its detector locus in the base\\
$N=2m!$ & degree of the full relative-Weyl splitting field\\
$d_L$ & absolute discriminant of that splitting field\\
$D_x$ & discriminant of the centre of the integral endomorphism ring
\end{tabular}
\end{center}
In Section~\ref{sec:polynomial-threshold} the adjoint component is
called $S$, and the original fine component is called $S_0$;
the finite map between them is always displayed.


\begin{theorem}[Polynomial discriminant forcing]\label{thm:forcing}
Assume the normalized adjoint level and the effective Hecke and Galois
inputs in Section~\ref{sec:polynomial-threshold}.  Let $x$ be a
relative-Weyl special point and let $L_x$ be the splitting field of
its Mumford--Tate torus.  Suppose $\operatorname{QRH}(\delta;L_x)$,
with $0<\delta\leq1/2$.  If
\begin{equation}\label{eq:forcing}
 \delta_W(x)\leq D,\qquad
 |\operatorname{disc}L_x|>
 \mathcal B_{S^{\rm ad},F,\delta}(H_D)^{1/\mu},
\end{equation}
then $\mathcal D_D=S$ and the Weil classes of every complex fibre
are algebraic.  The constants $\mu$ and $\mathcal B$ are given in
\eqref{eq:fixed-field-threshold}; for fixed datum, level, field and
$\delta$, the latter is a constant times $H_D^{2^{s-1}}$ times a
fixed power of a linear polynomial in $\log H_D$, where $s=\dim S$.
Cureton--Thorner supplies $\delta=1/32$, and Hu supplies $\delta=1/8$,
subject to their respective analytic theorems.  This implication
does not use a Galois-orbit bound in an endomorphism discriminant.
\end{theorem}

The meaning of the theorem is precise: finding one detector may
solve an entire family, but its degree and the discriminant of its
point must satisfy the \emph{same} inequality. Separate statements
that detectors exist at every CM point and that discriminants are
unbounded do not provide such a pair.

\begin{theorem}[The stronger threshold from CM heights]
\label{thm:sharper-forcing-intro}
For the principal unitary components of
Proposition~\ref{prop:ea-fine-normalization}, let the abelian
dimension be $m=4$ or $8$, put $N=2m!$, and let $x$ be full
relative-Weyl, meaning the natural $S_m\times C_2$ action specified
in the definition preceding Lemma~\ref{lem:weyl-signed-discriminant}. In addition to the Hecke hypotheses above, assume
the zero-free hypotheses of Theorem~\ref{thm:explicit-siegel-orbit}
for the CM endomorphism field of $A_x$.
Define the positive closed expression $\mathcal X(H)$ by
\eqref{eq:hybrid-X}. A detector of degree at most $D$ at $x$ forces
$\mathcal D_D=S$ whenever
\[
 |\operatorname{disc}L_x|>
 \bigl(e\mathcal X(H_D)\log\mathcal X(H_D)\bigr)^{mN/2}.
\]
The threshold is $H_D^{768+o(1)}$ for $m=4$ and
$H_D^{10\,569\,646\,080+o(1)}$ for $m=8$, with all coefficients
and logarithmic powers specified in Section~\ref{subsec:hybrid-point-forcing}.
For the centre-order discriminant the corresponding sufficient
threshold is
$\bigl(e\mathcal X(H_D)\log\mathcal X(H_D)\bigr)^{2m}$.
\end{theorem}

The proof combines the published average Colmez formula and the
explicit endomorphism-discriminant estimate of Gaudron--R\'emond
with the named zero-free input and Jaswal's Hecke construction.
At a fixed total degree one can take the smaller of this threshold
and the threshold in Theorem~\ref{thm:forcing}.
Corollary~\ref{cor:hybrid-detector} proves the statement. The cycle
degree $D$ and the degree $H_D$ of its parameter image remain
different quantities.

\begin{theorem}[One prime with an arithmetic certificate]
\label{thm:arithmetic}
Let $X$ be smooth and projective over
$R=\mathcal O_L[1/b]$, where $0\ne b\in\mathcal O_L$, with an ample line bundle, and let
$\mathscr H_D$ be a projective parameter space for pairs of
effective cycles of degree at most $D$, with universal cycle
classes. Fix affine charts $\Spec B_j$ covering $\mathscr H_D$,
and choose a nonzero $a_D\in\mathcal O_L$ that annihilates every
\begin{equation}
 \ker(B_j\longrightarrow B_j\otimes_R L).
 \label{eq:torsion}
\end{equation}
Such an $a_D$ exists. For the test it is enough to compute an
integer for which all localized chart rings are flat; explicit
chart equations and ideal-membership certificates give one.
A power of this integer annihilates the torsion ideals. At every good place
$v\nmid ba_D$, an equality
\begin{equation}
 M\operatorname{sp}_{v,\ell}(t)
   =\cl_\ell Z_v^+-\cl_\ell Z_v^-,\qquad
 \deg Z_v^\pm\leq D,\quad M\geq1,
 \label{eq:oneprime-eq}
\end{equation}
for a rational Betti class $t$, one geometric place and one
allowable $\ell$, lifts to an equality with the same $M$ and the
same degree bound in characteristic zero. The analogous statement
holds for a detector and its nonzero projected class.
\end{theorem}

This is an effective test once the model and the parameter-space
equations are given. It does not assert that those equations or
their certificate have been evaluated for an unspecified level or
an unspecified CM point.

\begin{theorem}[One relatively flat detector]\label{thm:flat}
Let $\mathcal S$ be integral and noetherian, and let a polarized abelian scheme
with its labelled quadratic action be defined over $\mathcal S$.
Fix a prime $\ell$ invertible there and a projective cycle
parameter space $\mathscr H_D\to\mathcal S$ with universal
classes. If $h\in\mathscr H_D$ is a geometric detector point
and this morphism is flat at $h$, then every geometric fibre has
a detector of degree at most $D$. In particular this holds on
every complex fibre. The point $h$ may have positive characteristic.
\end{theorem}

For the full Gaussian eightfold family no detector satisfying
either Theorem~\ref{thm:forcing} or Theorem~\ref{thm:flat} is
constructed here. These statements are conditional criteria for
Weil's Hodge classes, with a specified geometric existence problem.

\section{Proof overview}

The first difficulty is to preserve the degree bound while moving
a cycle. Proper Hilbert parameter spaces solve that problem:
universal cycle classes are locally constant, so the condition
that their Weil projection is nonzero selects whole connected
components. Their image in the base is closed. A relative-Weyl
point cannot lie on a proper positive-dimensional special
subvariety in adjoint type $A$. If the closed image were proper,
the point would be a maximal zero-dimensional special subvariety
of one of its components. The effective Andr\'e--Oort input then
bounds its discriminant. The general maximal-special theorem bounds the field of a geometric
component using total degree.  The stronger Hecke argument keeps the
base field fixed by taking the Galois union of that component.  This
is what makes the resulting degree dependence polynomial.

For a fixed arithmetic fibre the obstruction is torsion in the
coordinate rings of the cycle parameter space over its ring of
integers. Killing that torsion makes the parameter space flat.
Flatness supplies a characteristic-zero generization, and smooth
proper specialization transports the cycle equality unchanged.
Clearing denominators in a Gr\"obner-basis certificate makes this
argument executable. In a family of positive dimension, flatness
over the moduli base is stronger; its open-image property and
properness give the second one-prime implication.

Prime distribution only enters after the forbidden primes have
been determined. For an abelian splitting field one can force
complete splitting and avoid every forbidden prime by a single
arithmetic progression. This does not establish relative flatness
or a detector degree bound. The nonabelian Weyl group explains
why this abelian prime bound and the relative-Weyl discriminant
criterion have different ranges of applicability.

\section{Closed detector loci and discriminant forcing}

This section proves the geometric implication independently of
any quantitative construction of cycles.

\begin{lemma}[An algebraic Weil projector]\label{lem:projector}
Choose an integral endomorphism $a\in K$ for which $a/\bar a$
is not a root of unity. There is a polynomial $P\in\Z[T]$ such
that $P(a^*+\bar a^*)$ on middle cohomology vanishes off the Weil
plane and acts there by a nonzero integer. Its normalization
$e_W$ is an algebraic projector, defined over the field of the
labelled action and compatible with smooth proper specialization.
If an algebraic class has nonzero $e_W$-projection, the entire
Weil plane consists of algebraic classes.
\end{lemma}
\begin{proof}
Over $K$, middle cohomology has exterior bidegrees
$(j,2n-j)$, for $0\leq j\leq2n$. The operator in the statement
has eigenvalue
\[
 \lambda_j=a^j\bar a^{2n-j}+\bar a^ja^{2n-j}.
\]
Each $\lambda_j$ is a rational algebraic integer, and
$\lambda_j=\lambda_{2n-j}$. If two of
$\lambda_0,\ldots,\lambda_n$ coincide, put $r=a/\bar a$ and
divide their equality by a nonzero power of $a\bar a$. The
identity $u+u^{-1}=v+v^{-1}$ implies $u=v$ or $u=v^{-1}$,
and hence a nonzero power of $r$ is one. Thus these eigenvalues
are distinct. Take
\[
 P(T)=\prod_{j=1}^{n}(T-\lambda_j).
\]
The two extreme bidegrees are exactly the Weil plane, so division
by $P(\lambda_0)\ne0$ gives $e_W$. Pullback by an endomorphism
and linear combinations and compositions of such operators are
algebraic correspondences and commute with specialization.

Let $w=e_W\cl(Z)\ne0$. Both $w$ and $a^*w$ are algebraic.
The eigenvalues of $a^*$ on the Weil plane are $a^{2n}$ and
$\bar a^{2n}$. They are distinct nonrational conjugates: a
rational value of $a^{2n}$ would make their quotient one.
Consequently this operator has no rational eigenline. The two
rational vectors $w,a^*w$ are independent and span the plane.
\end{proof}

\begin{lemma}[Closedness and fields of components]\label{lem:closed}
The locus $\mathcal D_D$ is closed and defined over $F$. If
\eqref{eq:total-degree} holds, every geometric component is
defined over an extension of $F$ of degree at most $H_D$.
\end{lemma}
\begin{proof}
Use the finite union of products of the full projective Hilbert
schemes for the finitely many Hilbert polynomials of integral
subvarieties of bounded degree. Include all their flat limits, and
allow repetition to represent multiplicities. Boundedness of their Hilbert polynomials is
\cite[Expos\'e~XIII, Corollary~6.11, pp.~658--660]{SGA6};
projectivity is \cite[Expos\'e~221, Theorem~3.2, p.~260]{FGA}. One uses a fixed very ample multiple of the
polarization and translates its degrees back to $\theta$.
The top-dimensional fundamental cycles in the universal flat
families give sections of the pulled-back lisse cohomology sheaf;
trace and smooth proper duality give this construction and its
base-change compatibility, as in \cite[Cycle, \S2.3]{SGA45}.

Apply $e_W$ to the universal section. On a connected parameter
component its value is either everywhere zero or everywhere
nonzero. Hence the nonzero locus is open and closed, and is
proper over $S$. Its image is closed. If an effective cycle detects, at least one of its irreducible
components detects and has no larger degree. Conversely every
irreducible detector occurs in the parameter space. Thus this
image is exactly $\mathcal D_D$. The action, projector
and class construction are defined over $F$, which proves
descent of the nonzero locus and its image.

The Galois group permutes the geometric components, whose number
is at most $H_D$. The stabilizer of a component has index at most
$H_D$. Galois descent of its defining ideal over the fixed field
of that stabilizer gives the asserted field of definition.
\end{proof}

\begin{lemma}[Maximal special points in type $A$]\label{lem:weyl}
A relative-Weyl special point in the stated datum lies on no
proper positive-dimensional special subvariety of $S$.
\end{lemma}
\begin{proof}
The image in the adjoint group of the Mumford--Tate group of a
special subvariety through the point contains its maximal torus.
After algebraic closure it is therefore a reductive subgroup of
full rank. Its roots form a subset of the ambient roots, stable
under the Galois action on that torus. This action contains the
full Weyl group. In each absolute type-$A$ factor the Weyl group
is transitive on the roots. A nonempty root subset therefore
contains every root of that factor. Rational simplicity makes
Galois transitive on the absolute factors. Thus the root subset
is either empty or the whole ambient root system. In the first
case the group is a torus and its domain has dimension zero. In
the second case it has the full adjoint group and the same
connected domain as $S$. This is the full-rank root argument
of Chai--Oort \cite[Lemma 6.3, pp.~621--622]{ChaiOort}.
\end{proof}

\begin{proposition}[The general maximal-special bound]\label{prop:general-forcing}
Assume Jaswal's effective maximal-special-subvariety theorem, and
let $x$ be a relative-Weyl special point. Write $L_x^{\rm ad}$ for
the splitting field of its adjoint Mumford--Tate torus. Suppose
that $\zeta_{L_x^{\rm ad}}(s)$ has no zero in
$\operatorname{Re}s>1-\delta$, apart from its pole at $1$.
If
\[
 \delta_W(x)\leq D,\qquad
 |\operatorname{disc}L_x^{\rm ad}|
   >c_S(\delta,H_D,[F:\Q]H_D),
\]
then $\mathcal D_D=S$, and the Weil classes of every complex fibre
are algebraic. Cureton--Thorner gives the indicated implication
with $\delta=1/32$; Hu gives it with $\delta=1/8$, subject to their
respective analytic theorems. No additional endomorphism-discriminant orbit bound is
needed for this implication.
\end{proposition}

\begin{proof}[Proof of Proposition~\ref{prop:general-forcing}]
Suppose $\mathcal D_D\ne S$, and choose a geometric component
$Z$ containing $x$. By Lemma~\ref{lem:closed}, its degree is at
most $H_D$ and it is defined over a number field of degree at
most $[F:\Q]H_D$. Lemma~\ref{lem:weyl} shows that $\{x\}$
is a maximal special subvariety contained in $Z$. Jaswal's
Theorem~1.1(ii) applies, including when $\dim Z=0$, to the
connected center of the image of its Mumford--Tate group in the
adjoint group. For a special point this is its adjoint torus.
It gives
\[
 |\operatorname{disc}L_x^{\rm ad}|\leq
 c_S(\delta,H_D,[F:\Q]H_D),
\]
contrary to the assumed discriminant inequality. Thus $\mathcal D_D=S$.
Lemma~\ref{lem:projector} gives algebraicity of each Weil plane.
\end{proof}

\section{A polynomial discriminant threshold after Chai--Oort and Jaswal}
\label{sec:polynomial-threshold}

The purpose of this section is to keep the field of definition fixed when
applying Jaswal's Hecke theorem to a detector locus.  His theorem applies
to a locus irreducible over that field; geometric irreducibility is not
required \cite[Theorem~6.7]{Jaswal}.  Chai--Oort's root calculation makes
this stronger form available for relative-Weyl points in a full unitary
family.

Let $S$ be a geometrically connected component, defined over a number
field $F$, of the Shimura variety of a connected adjoint
$\mathbb Q$-simple group $G$ of absolute type $A$.  Assume that $F$
contains the reflex field and that the chosen neat product level
satisfies Jaswal's normalization (Lv): at $3$ it is contained in the
principal congruence subgroup for the chosen faithful representation.
Write $s=\dim S\geq1$ and $d=[F:\mathbb Q]$.  Degree means degree for
the Baily--Borel line bundle.  For a reduced closed locus, its
\emph{total degree} is the sum of the degrees of all geometric
irreducible components, including components of smaller dimension.


\subsection{The threshold and its constants}

We use the fixed-datum constants in \cite[Table~1 and \S2.3]{Jaswal}.
Thus $n_T$ bounds the splitting-field degrees of the relevant tori,
$R_G$ bounds the degrees of the Galois closures of their reflex
fields, $a$ bounds the auxiliary level indices, and $r_U$ bounds the
numbers of cyclic factors of the local torus-level quotients.
The positive constants $B,c_U,k_1,t_\rho,s_\rho,C_0$ have the
meanings specified there.  In particular,
\[
 f=(\dim G+1)^2,\qquad
 \mu=\frac{1}{16t_\rho s_\rho n_T(n_T-1)},\qquad B'=Bc_U.
\]
They are fixed before the detector degree or the point is chosen.
They include the chosen level and representation; a dimension alone
does not determine all of them.  Put
\[
 e_s=\left\lfloor\sum_{k=1}^{s}k2^{1/k}\right\rfloor+2,
 \qquad c_B=e_s^{2s}.
\]
For $0<\delta\leq\tfrac12$ let
\[
 \kappa_\delta=\frac{2-\delta}{\delta(1+\delta)},\qquad
 \vartheta=2\mu,
\]
and define
\begin{align*}
 \beta_0={}&\max\left\{
 \frac{128\kappa_\delta}{\vartheta^2\delta^2},\,
 \sqrt{\frac{96}{\vartheta^3}},\,
 3n_T\log n_T+1,\,\frac{64n_T}{\vartheta}\right\},\\
 C_*={}&\frac{\exp(-\mu\beta_0)}{R_Ga(a!)^{r_U}},\\
 \sigma={}&\max\left\{
 \frac{64\kappa_\delta+1}{\mu}+8n_T,\,
 1+48n_T+4C_0+8n_T(e^2/B'+2)\right\}.
\end{align*}
These are Jaswal's constants for the prime and Galois estimates
\cite[\S3.2, Corollary~3.4, and Lemma~6.8]{Jaswal}.
Set
\begin{align*}
 q&=2^{s-1},& b&=(k_1d+3f)q,&
 \Lambda(H)&=\frac{(c_BdH)^q}{C_*},\\
 W(\Lambda,b)&=2\log\Lambda+1+
       \frac{2b}{\delta}\log\frac{2\sigma b}{e\delta}.
\end{align*}
The threshold used below is
\begin{equation}\label{eq:fixed-field-threshold}
 \mathcal B_{S,F,\delta}(H)
 =\Lambda(H)\left[\sigma\bigl(W(\Lambda(H),b)+1\bigr)\right]^{b/\delta}.
\end{equation}

For a special point $y$, let $L_y$ be the splitting field of its
adjoint Mumford--Tate torus, let $d_{L_y}$ be its absolute field
discriminant, and let $\alpha_y\geq1$ be the Klingler--Yafaev level
invariant with the normalization of \cite[Definition~2.2]{Jaswal}.
Write $\operatorname{QRH}(\delta;L_y)$ for
$\zeta_{L_y}(z)\ne0$ when $\operatorname{Re}z>1-\delta$ and $z\ne1$.

\begin{theorem}[Fixed-field form of the Hecke bound]
\label{thm:fixed-field-bound}
Assume Jaswal's Corollary~3.4, Theorem~6.7, and Lemma~6.8 with the
constants above.  Let $Y\subsetneq S$ be a reduced closed
$F$-subvariety of total degree at most $H\geq1$.  If $y\in Y$ is
relative Weyl and $\operatorname{QRH}(\delta;L_y)$ holds, then
\[
 \alpha_y d_{L_y}^{\mu}\leq\mathcal B_{S,F,\delta}(H).
\]
In particular,
\[
 d_{L_y}\leq\mathcal B_{S,F,\delta}(H)^{1/\mu},
 \qquad \alpha_y\leq\mathcal B_{S,F,\delta}(H).
\]
\end{theorem}

\begin{proof}
Choose a geometric irreducible component $C$ of $Y$ containing $y$.
Let $Z$ be the reduced union of its
$\operatorname{Gal}(\overline F/F)$-conjugates.  It is a closed
$F$-irreducible subvariety with total degree at most $H$.  Write
$z=\dim Z\leq s-1$ and
$\Omega_y=C_*\alpha_y d_{L_y}^{\mu}$.

First suppose $z=0$.  The $F$-orbit of $y$ is contained in $Z$ and
has at most $H$ points.  If $E_y$ is the reflex field of its torus
datum, Jaswal's Galois bound gives at least $\Omega_y$ conjugates
over $E_y$ \cite[Corollary~3.4 and Lemma~6.4]{Jaswal}.
Passing from $E_y$ to $FE_y$ loses a factor at most
$[FE_y:E_y]\leq d$.  Its orbit over $F$ is at least as large as
its orbit over $FE_y$, and hence $\Omega_y\leq dH$.
Since $c_B\geq1$, $q\geq1$, $C_*\leq1$, and the bracket in
\eqref{eq:fixed-field-threshold} is at least $1$, the asserted
inequality follows.

Now suppose $z>0$.  The special closure of $C$ contains $y$ and is
positive dimensional.  Lemma~\ref{lem:weyl} makes it all
of $S$, so $C$ is Hodge generic.  Every conjugate of $C$ is also
Hodge generic, since Galois conjugates of special subvarieties are
special.  The point $y$ is not strongly special in the adjoint
variety: its maximal torus is nontrivial.

Apply \cite[Theorem~6.7]{Jaswal} with $P=G$, $V=\{y\}$ and
$\bar d=d$.  A larger special subvariety of $Z$ through $y$ is
impossible by Lemma~\ref{lem:weyl}, since $Z$ is proper
in $S$.  Therefore that theorem's strict inequality fails.  In its
notation this says
\begin{align*}
 C_*\alpha_y d_{L_y}^{\mu}
 &\leq (c_Bd\deg Z)^{2^z}
       x_y^{(k_1d+3f)2^z}\\
 &\leq (c_BdH)^q x_y^b.
\end{align*}
Here $x_y\geq1$ is the explicit split-prime threshold of
\cite[\S6.2]{Jaswal}.  Divide by $C_*$ and use
\cite[Lemma~6.8]{Jaswal}; its function $\mathcal B$ is exactly
\eqref{eq:fixed-field-threshold}.  The remaining two bounds follow
from $\alpha_y\geq1$ and $d_{L_y}\geq1$.
\end{proof}

The field in this argument is $F$ throughout.  Taking the Galois
union of a geometric component is precisely what the
$F$-irreducible version of Jaswal's theorem permits.

\begin{corollary}[Explicit polynomial form]\label{cor:pure-polynomial}
Put
\[
 C=\frac{(c_Bd)^q}{C_*},\qquad a_0=\frac b\delta,\qquad
 u=2\log C+2+2a_0\log(2\sigma a_0/e).
\]
For every $\varepsilon>0$ and $H\geq1$,
\[
 \mathcal B_{S,F,\delta}(H)
 \leq C\left[\sigma\left(u+\frac{2qa_0}{e\varepsilon}\right)\right]^{a_0}
 H^{q+\varepsilon}.
\]
Thus the discriminant exponent may be any number strictly greater
than $2^{s-1}/\mu$, with the displayed coefficient.
\end{corollary}

\begin{proof}
The definition gives
\[
 \mathcal B_{S,F,\delta}(H)
 =CH^q\left[\sigma(u+2q\log H)\right]^{a_0}.
\]
For $t\geq0$, the maximum of $t\exp(-\varepsilon t/a_0)$ is
$a_0/(e\varepsilon)$.  Set $t=\log H$.  Since
$H^{\varepsilon/a_0}\geq1$ and $u>0$,
\[
 u+2q\log H\leq
 \left(u+\frac{2qa_0}{e\varepsilon}\right)
 H^{\varepsilon/a_0}.
\]
Raising this inequality to the power $a_0$ proves the assertion.
\end{proof}

\subsection{Application to a fine unitary family}

Let $S_0$ be the fine unitary component carrying the abelian scheme.
Choose the fixed level so that its adjoint map $f:S_0\to S$ has
the normalized level used above.  The map is finite and surjective
on a connected component and satisfies
$f^*\mathcal L_S=\mathcal L_{S_0}$
\cite[Proposition~9.4]{Jaswal}.  For any closed reduced locus $D$
the projection formula, applied component by component, gives
\[
 \operatorname{tdeg} f(D)\leq\operatorname{tdeg}D.
\]
Repeated image components are counted only once on the left.  If
$D\ne S_0$, then $\dim D<\dim S_0$, so the finiteness of $f$
also gives $f(D)\ne S$.

\begin{corollary}[A detector forcing criterion]
\label{cor:polynomial-detector}
Assume the hypotheses of Theorem~\ref{thm:fixed-field-bound} at the
adjoint level.  Let $\mathcal D_D\subset S_0$ be the closed
$F$-locus of fibres admitting an effective detector of degree at
most $D$.  Suppose that its total Baily--Borel degree is at most
$H(D)$.  If a relative-Weyl point $x\in\mathcal D_D$ satisfies
$\operatorname{QRH}(\delta;L_x)$ and
\[
 d_{L_x}>\mathcal B_{S,F,\delta}(H(D))^{1/\mu},
\]
then $\mathcal D_D=S_0$.
\end{corollary}

\begin{proof}
If $\mathcal D_D$ were proper, its adjoint image would be a proper
closed $F$-locus of total degree at most $H(D)$.  The image of $x$
is relative Weyl with the same adjoint torus.  Theorem
\ref{thm:fixed-field-bound} would give the opposite discriminant
inequality.
\end{proof}

Theorem~\ref{thm:explicit-theta-chow-bound} supplies an explicit
$H(D)$ for the principal unitary families.  This criterion does not construct a
detector.  Its polynomial is in the degree of the locus in the base,
not in the degree of a cycle in a fibre.

\subsection{The splitting field in the balanced unitary case}

\begin{lemma}\label{lem:balanced-splitting-field}
Let $x$ be a CM point in a family of Weil type over an imaginary
quadratic field $K$, of $K$-dimension $2n$ and signature $(n,n)$.
The full Mumford--Tate torus of $x$ and its image in the adjoint
unitary group have the same splitting field.
\end{lemma}

\begin{proof}
Let $\nu$ be the similitude character.  On either complex
$K$-eigenspace the Hodge homomorphism $h$ has $n$ eigenvalues of
each Hodge type.  Hence
$\det_Kh(z)=\nu(h(z))^n$.  By the definition of the
Mumford--Tate group it is contained in
\[
 G^\sharp=\{g\in\operatorname{GU}(V):\det_Kg=\nu(g)^n\}.
\]
A scalar $a\in K^\times$ lies in this group precisely when
$(a/\bar a)^n=1$.  The identity component of this scalar group is
the split torus $\mathbb G_m$.  It contains the weight torus.
For the Mumford--Tate torus $T_x$ the kernel of
$T_x\to T_x^{\mathrm{ad}}$ therefore has split connected part.

Tensor the character sequence with $\mathbb Q$.  Its quotient is
a trivial Galois representation; possible finite components of
the kernel disappear after this tensor product.  The action factors
through a finite Galois group and is semisimple in characteristic
zero.  Consequently the rational character representation of $T_x$
is the direct sum of that of $T_x^{\mathrm{ad}}$ and a trivial
representation.  The two actions have the same kernel.  Rationalizing
an integral character lattice does not change the kernel of its
Galois action, so the splitting fields agree.
\end{proof}

\begin{proposition}\label{prop:no-abelian-weyl}
For the signature-$(n,n)$ unitary datum with $n\geq2$, a
relative-Weyl special point has nonabelian torus splitting field.
For $n=2$ and $n=4$ the splitting-field degrees are respectively
$48$ and $80640$.
\end{proposition}

\begin{proof}
The absolute root system is $A_{2n-1}$ and its Weyl group is
$S_{2n}$.  The Galois group of the splitting field acts faithfully
on the character lattice and contains this Weyl group.  It is
therefore nonabelian for $n\geq2$.

For the degree assertion put $m=2n$.  Since the group is adjoint,
its roots span its character space and the splitting group embeds
in $\operatorname{Aut}(A_{m-1})$.  The projection to the diagram
automorphism group is the nontrivial quadratic character of the
unitary field $K$.  At a relative-Weyl point the image contains
$S_m$ and surjects onto this diagram group of order $2$, so it is
all of $\operatorname{Aut}(A_{m-1})$, of order $2m!$.
Lemma~\ref{lem:balanced-splitting-field} identifies this with the
full Mumford--Tate splitting field.  Substitute $m=4$ and $m=8$.
\end{proof}

Alp\"oge's Dirichlet theorem gives the half-plane
$\operatorname{Re}z>\tfrac12$ for an abelian splitting field,
because its Dedekind zeta function is a product of Dirichlet
$L$-functions.  Proposition~\ref{prop:no-abelian-weyl} shows that
this improvement does not apply to the relative-Weyl points used
in Corollary~\ref{cor:polynomial-detector} for either of these
families.  It remains available for a one-prime argument based on
an independent geometric deformation certificate at a non-Weyl CM
point.

For the two base dimensions the constants depending only on the
dimension are
\[
\begin{array}{c|c|c|c|c}
\text{signature}&s&e_s&f&q\\ \hline
(2,2)&4&15&256&8\\
(4,4)&16&150&4096&32768
\end{array}
\]
and $c_B=15^8$ or $150^{32}$.  The values of the analytic
parameter enter through
\[
\begin{array}{c|c|c}
\delta&\kappa_\delta&b/\delta\\ \hline
1/32&672/11&32b\\
1/8&40/3&8b\\
1/2&2&2b
\end{array}
\]
The last row requires the corresponding Riemann hypothesis for
the particular nonabelian field if it is to be used at a
relative-Weyl point.

\begin{proof}[Proof of Theorem~\ref{thm:forcing}]
Corollary~\ref{cor:polynomial-detector} applies to the adjoint image
of the fine family.  Lemma~\ref{lem:balanced-splitting-field}
identifies the splitting fields used in the two statements.
Lemma~\ref{lem:projector} then extracts the entire Weil plane.
\end{proof}

\section{A fixed fine level and numerical unitary constants}
\label{sec:ea-numerical-unitary}

This section fixes the group, integral representation and levels used in
our effective criterion.  It gives closed numerical expressions for the
fixed-datum constants, including the field of one fine component and a
bound for the degree of its adjoint map.
An explicit choice of $H(D)$ for the principal unitary families
is given in Theorem~\ref{thm:explicit-theta-chow-bound}.

\subsection{The fine component and its adjoint map}

Let $K=\mathbf Q(i)$ or $\mathbf Q(\omega)$, and let
$m=2n\in\{4,8\}$.  On $K^m$ put
\[
 J=\begin{pmatrix}0&I_n\\I_n&0\end{pmatrix},\qquad
 h(v,w)=\overline v^{\,t}Jw,
\]
and define the connected reductive group
\[
 G^\sharp=\{g\in\operatorname{GU}(h):
                   \det_K(g)=\nu(g)^n\}.
\]
Its derived group is $\operatorname{SU}(h)$ and its adjoint group is
$G=\operatorname{PU}(h)$.  The multiplier gives an exact sequence
\begin{equation}\label{eq:ea-sharp-exact}
 1\longrightarrow\operatorname{SU}(h)\longrightarrow G^\sharp
 \xrightarrow{\nu}\mathbf G_m\longrightarrow1,
 \qquad
 a\longmapsto\operatorname{diag}(aI_n,I_n)
\end{equation}
with the displayed rational section.

Choose an imaginary generator $\delta$ of the inverse different of
$K/\mathbf Q$: one may take $\delta=(2i)^{-1}$ or
$\delta=(1+2\omega)^{-1}$.  The alternating form
\[
             E(v,w)=\operatorname{Tr}_{K/\mathbf Q}(\delta h(v,w))
\]
is integral and perfect on $\Lambda_0=\mathcal O_K^m$.  Indeed the
trace dual of $\mathcal O_K$ is its inverse different and $J$ is
unimodular.  Thus $G^\sharp$ embeds in the symplectic similitude group
of a principal symplectic lattice of rank $2m$.
Let $K^\sharp_p$ be its integral stabilizer for $p\ne3$, and put
\[
 K^\sharp_3=G^\sharp(\mathbf Q_3)\cap
       \ker\bigl(\operatorname{GL}(\Lambda_0\otimes\mathbf Z_3)
             \longrightarrow\operatorname{GL}(\Lambda_0/3\Lambda_0)\bigr).
\]
Write $K^\sharp=\prod_pK^\sharp_p$.  This is a neat product level:
eigenvalues of a matrix congruent to $1$ modulo $3$ lie in
$1+3\overline{\mathbf Z}_3$, and the multiplicative group they generate
contains no nontrivial root of unity.  The same statement holds after
conjugating the level.

\begin{proposition}\label{prop:ea-fine-normalization}
For the balanced weight-one Shimura datum of $G^\sharp$, the reflex
field is $\mathbf Q$.  Its level-$K^\sharp$ Shimura variety has two
geometric connected components, each defined over
\[
                         F=\mathbf Q(\zeta_3),\qquad [F:\mathbf Q]=2.
\]
Let $S_0$ be either component.  There is a choice of canonical adjoint
level and a finite \emph{\'etale} $F$-map $f:S_0\to S$ onto an adjoint
connected component for which
\begin{equation}\label{eq:ea-hf-values}
 \deg f\le h_f,\qquad
 h_f=\begin{cases}64,&m=4,\\32768,&m=8.\end{cases}
\end{equation}
For this adjoint level, the cutoff $C_0=4$ is valid in Jaswal's
Table~1.  The fine component carries the pullback of the universal
principally polarized abelian scheme with level $3$.
\end{proposition}

\begin{proof}
Over an algebraic closure, the Hodge cocharacter on the standard
$K$-embedding block has two eigenvalue multiplicities $n,n$, and its
multiplier is $z$.  Its conjugacy class is fixed by the nontrivial
embedding of $K$, which exchanges these two multiplicities.  Thus the
reflex field is $\mathbf Q$.  Equivalently, the rational section in
\eqref{eq:ea-sharp-exact}, as a cocharacter, represents that complex
conjugacy class.  Its image in the abelianization has cocharacter $z$.

The multiplier image of the level is
\[
 \nu(K^\sharp)=
 (1+3\mathbf Z_3)\times\prod_{p\ne3}\mathbf Z_p^\times.
\]
The inclusion follows from integrality and, at $3$, from the perfect
symplectic form.  Equality follows from the section in
\eqref{eq:ea-sharp-exact}.  The derived group is simply connected and
has a noncompact real group.  Strong approximation therefore identifies
the connected-component set with
\[
 \mathbf Q_{>0}^{\times}\backslash\mathbf A_f^\times/
                    \nu(K^\sharp)\simeq(\mathbf Z/3\mathbf Z)^\times.
\]
One can verify the double-coset assertion directly: match multipliers
using the rational section and an element of the level; the remaining
adelic element lies in $\operatorname{SU}(h)(\mathbf A_f)$ and is
removed by strong approximation.  The canonical Galois action on this
quotient is the cyclotomic character modulo $3$, since the multiplier
cocharacter is $z$.  This is the connected-component reciprocity law
of \cite[Theorem~5.17 and \S13, formula~(64), pp.~59 and~119]{MilneShimura}.
It is also the action on the two components of the Siegel level-$3$
moduli space, as described in \cite[p.~125]{MilneShimura}; the natural
map induces the same multiplier quotient.  Each component is consequently
defined over $\mathbf Q(\zeta_3)$.  Pulling back the universal abelian
scheme from the fine Siegel level gives the last assertion.

For the degree bound the adjoint level must be specified carefully.
The natural integral model of $G^\sharp$ is reductive over
$\mathbf Z[1/6]$: after the quadratic \emph{\'etale} splitting it is
\[
       \{(A,t)\in\operatorname{GL}_m\times\mathbf G_m:
                                      \det A=t^n\},
\]
a connected reductive group with derived group $\operatorname{SL}_m$.
The standard representation is a closed immersion, so its schematic
closure agrees with this model at every $p>3$.
Take $\mathcal S=\{2,3\}$ and use precisely Jaswal's adjoint level
\cite[\S9.1, Lemma~9.2]{Jaswal}:
\begin{equation}\label{eq:ea-canonical-adjoint-level}
 K^{\rm ad}_p=\begin{cases}
       \lambda(K^\sharp_p),&p=2,3,\\
       \mathcal G^{\rm ad}(\mathbf Z_p),&p>3.
                      \end{cases}
\end{equation}
The Lie lattice
$\Lambda_{\rm ad}=\operatorname{Lie}(G^{\sharp,\rm der})
                   \cap\operatorname{End}_{\mathbf Z}(\Lambda_0)$
is preserved by this level.  At $3$ its adjoint matrices are congruent
to $1$ modulo $3$, because conjugation by $1+3M$ changes every integral
Lie matrix by a multiple of $3$.  At $p>3$ the level is hyperspecial.
Thus it is a neat product level and $C_0=4$ works.  No enlargement of
\eqref{eq:ea-canonical-adjoint-level} at $2$ or $3$ is implicit: such an
enlargement would need an additional finite level-index factor in
\eqref{eq:ea-hf-values}.

Here are all the centre parameters in Jaswal's degree bound.  Over $K$
a central scalar has coordinates $(a,b)$ and satisfies $(a/b)^n=1$.
Consequently
\[
 X^*(Z(G^\sharp))=
      (\mathbf Ze_1\oplus\mathbf Ze_2)/\mathbf Z n(e_1-e_2).
\]
Conjugation exchanges $e_1,e_2$.  The free quotient has rank one and
trivial action; the torsion is $\mathbf Z/n\mathbf Z$, with action
$-1$.  In the notation of \cite[\S9.3]{Jaswal},
\[
 r_Z=1,\quad m_Z=n,\quad q=1,\quad L_1=\mathbf Q,\qquad
 L_0=\begin{cases}\mathbf Q,&n=2,\\K,&n=4.\end{cases}
\]
In particular $L_0$ is \emph{not} $K$ in the $n=2$ case: inversion
on $\mathbf Z/2\mathbf Z$ is trivial.  These fields are the splitting
fields of the character modules, not necessarily the fields containing
all geometric points of the corresponding finite group.
Both possible $L_0$ have class number one.  For the two quadratic
fields this follows from their norm-Euclidean rings of integers.
The connected centre is the scalar $\mathbf G_m$.  Its maximal compact
subgroup has index one at $2$ and index two at $3$ in its intersection
with the fine level.  Hence $\iota_{K^\sharp}=2$.

The integral-model hypotheses used in that bound hold here at every
$p>3$.  The centre is the displayed group of multiplicative type;
its maximal subtorus is scalar $\mathbf G_m$, with finite flat quotient
having the displayed torsion character module.  These modules are
unramified outside $2,3$.  The quotient becomes $\mu_n$ over the
unramified splitting algebra.  Lang's theorem and Hensel's lemma give
the required vanishing for integral torus torsors; finiteness gives
extension of generic points over a discrete valuation ring; Kummer
theory gives the unit description of $\mu_n$-torsors.  Finally the
integral central points are exactly the central intersections with the
integral stabilizer.  These verify conditions (M1)--(M6) in \cite[\S9.3]{Jaswal}.

Corollary~9.10 of that source therefore gives
\[
 h_f\le
 2n^{\log_2 [L_0:\mathbf Q]+[L_0:\mathbf Q]+4}
 =\begin{cases}2\cdot2^5=64,&n=2,\\
                 2\cdot4^7=32768,&n=4.
   \end{cases}
\]
The map and its finite \emph{\'etale} property on components are
\cite[Proposition~9.4]{Jaswal}; its canonical construction is over
$\mathbf Q$, and restriction of a component and its image is over $F$.
\end{proof}

A universal theta characteristic is not needed for a relative very
ample line bundle.  If $\mathcal P$ is the normalized Poincar\'e bundle
and $\lambda:\mathcal A\to\mathcal A^\vee$ the principal polarization,
then $(1,\lambda)^*\mathcal P$ has fibre class $2\theta$.
Its third power has class $6\theta$ and is relatively very ample, by the
very-ampleness of the third power of an ample line bundle on an abelian
variety \cite[Remark~7.3, p.~12]{MilneAV}, together with cohomology
and base change.
This fixes a fibre polarization convention; it does not supply equations
for a Baily--Borel model or a bound for a detector's image degree.

\subsection{Closed formulas for the datum constants}

The following bounds apply at the adjoint level
\eqref{eq:ea-canonical-adjoint-level}.  They also apply to the full
integral adjoint level with principal congruence at $3$, but the degree
bound \eqref{eq:ea-hf-values} refers to the level just specified.
Set
\begin{equation}\label{eq:ea-integer-parameters}
 \begin{aligned}
 r&=m-1,& d_\rho&=m^2-1,& N&=2m!,\\
 M&=2^r,& C&=(rM)^r,\\
 A&=(NC)!\,M\,(M!)^4,&
 B_\rho&=2r!\,M C^{r-1}+1,\\
 C_\Phi&=(2d_\rho rM)^r,&
 a&=\prod_{j=0}^{d_\rho-1}(3^{d_\rho}-3^j).
 \end{aligned}
\end{equation}
These are explicit integers, however large their decimal expansions.
In particular no unspecified finite search occurs in these definitions.

\begin{proposition}\label{prop:ea-explicit-constants}
Valid uniform choices in Jaswal's Table~1 are
\begin{equation}\label{eq:ea-complete-constants}
 \begin{gathered}
 n_T=R_G=N,\qquad t_\rho=M,\qquad s_\rho=d_\rho,
 \qquad r_U=2r,\qquad C_0=4,\\
 B=(A\,a!)^{-2r},\qquad c_U=(2C_\Phi)^{-1},\qquad B'=Bc_U,\\
 k_1=2A d_\rho^3B_\rho,\qquad
 f=(d_\rho+1)^2=m^4,\qquad
 \mu_{\rm unif}=\frac1{16Md_\rho N(N-1)}.
 \end{gathered}
\end{equation}
For the relative-Weyl special-point application, the orbit lower-bound
proof may instead use
\begin{equation}\label{eq:ea-point-mu}
 t_\rho=1,\qquad s_\rho=m(m-1),\qquad
 \mu=\frac1{16m(m-1)N(N-1)},
\end{equation}
while retaining the other conservative constants in
\eqref{eq:ea-complete-constants}.
\end{proposition}

\begin{proof}
We give the arithmetic parameters needed by the arguments of Ullmo--Yafaev and Edixhoven--Yafaev.  Their geometric and reciprocity assertions remain the
named source inputs used in Jaswal's theorem.

\emph{Field degrees and central isogenies.}
A maximal $\mathbf Q$-torus of the adjoint group has the $A_{m-1}$
root lattice as character lattice.  Its splitting Galois group embeds
in $\operatorname{Aut}(A_{m-1})$, of order $2m!=N$.
Every torus in $G$ is contained in a maximal $\mathbf Q$-torus, giving
$n_T=N$.  For a Shimura subgroup $H$, extend a maximal
$\mathbf Q$-torus of $H$ to one of $G$.  The latter's Galois splitting
field splits $H$ and contains its reflex field, hence also its reflex
field's Galois closure.  This gives $R_G=N$.

For an irreducible root system of rank $s$, the order of the centre of
its simply connected group is at most $2^s$: the values are $s+1$ in
type $A_s$, $2$ in types $B_s,C_s$, $4$ in type $D_s$, and $3,2,1,1,1$
in types $E_6,E_7,E_8,F_4,G_2$.  Products multiply these orders and add
the ranks.  Consequently both
\[
 |Z(H)^\circ\cap H^{\rm der}|\le M
 \quad\hbox{and}\quad
 \deg(H^{\rm sc}\to H^{\rm der})\le M.
\]
These are the two uniform bounds used in
\cite[Lemma~2.4 and the proof of Lemma~2.8, pp.~11 and~17--19]{UllmoYafaev}.

\emph{Hodge and reciprocity coordinates.}
Put $T=Z(H)^\circ$, $C_H=H/H^{\rm der}$, and let
$m_0\le M$ be the degree of $\alpha:T\to C_H$.
Conjugate the Hodge cocharacter $\mu_H$ into a maximal torus of $H$.
Its central projection is its average under the Weyl group of
$H^{\rm der}$.  That Weyl group permutes the weights of
$\operatorname{Lie}G$, all of whose pairings with $\mu_H$ belong to
$\{-1,0,1\}$.  Each central-weight pairing with this average therefore
has absolute value at most one.  The representation
$t\mapsto\rho(t)^{m_0}$ descends to $C_H$, with at most $d_\rho$
weights and with Hodge weights at most $m_0\le M$.
Thus $s_\rho=d_\rho$ and $t_\rho=M$ work.

Let $F_H$ be the Galois closure of the reflex field and use the
reciprocity map
$r_C:\operatorname{Res}_{F_H/\mathbf Q}\mathbf G_m\to C_H$.
The map $r_T$ defined by $\alpha r_T=r_C^{m_0}$ identifies $X^*(T)$
with a sublattice of $\mathbf Z^{[F_H:\mathbf Q]}$.
In these coordinates the weights of $\rho|_T$, and the corresponding
descended characters of $C_H$, have entries of absolute value at most
$M$.  This is the coordinate calculation in
\cite[Lemma~2.6, pp.~12--14]{UllmoYafaev}, with the preceding Weyl average
making its bound explicit.  The weights of $T$ generate $X^*(T)$
integrally, since $\rho|_T$ is faithful.  The descended characters span
$X^*(C_H)$ rationally.

\emph{A common reciprocity exponent.}
Let $L_C$ be the image of $X^*(C_H)$ in
$\mathbf Z^{[F_H:\mathbf Q]}$, of rank $t\le r$.  It contains $t$
independent vectors with all coordinates bounded by $M$.  A nonzero
$t$-row coordinate minor has determinant at most
$t!M^t\le(rM)^r=C$.  The torsion order of
$\mathbf Z^{[F_H:\mathbf Q]}/L_C$ divides such a minor, and is therefore
at most $C$.

In \cite[Lemma~2.7, pp.~15--17]{UllmoYafaev}, the product of the nonzero Smith
invariants is exactly this torsion order.  A section up to that integer,
averaged over the Galois group of $F_H$, gives a rational morphism whose
composition with reciprocity is a power of exponent at most $NC$.
Thus $(NC)!$ is a common multiple of every exponent required there,
and of $2$.  It is a common \emph{exponent}, not merely a larger number.
In \cite[Lemma~2.8, pp.~17--19]{UllmoYafaev} take $n_2=n_3=M!$ to kill the
two central-isogeny cokernels.  The preliminary real component factor
is killed by $M$, since a real torus has elementary $2$ component group
of order at most $2^r=M$.  The remaining exponent in that proof is
$n_3n_2^3$.  Consequently $M(M!)^4$ is valid there, and
\cite[Proposition~2.9, pp.~19--20]{UllmoYafaev} takes the product of the two
exponents.  This proves the stated value of $A$.

\emph{The EY coordinate bound.}
Let $L_T$ denote the image of $X^*(T)$ in the same integral coordinate
space.  Select $t$ coordinate positions that give an injective projection
$L_T\to\mathbf Z^t$.  Its image is generated by projected weights of
size at most $M$, so its index in $\mathbf Z^t$ is at most $C$.
A Hermite-normal-form basis has entries at most that index.  Cramer's
rule then bounds the coordinates of each weight in that basis by
\[
                    t!M C^{t-1}\le r!M C^{r-1}.
\]
Differences of weights are bounded by twice this integer.  The added
$1$ in $B_\rho$ also covers a strict-inequality formulation.
This is exactly the difference-of-weight bound needed in
\cite[Proposition~7.4.3, author version p.~23]{EdixhovenYafaev}.
Jaswal's formula gives $k_1=2Ad_\rho^3B_\rho$.

\emph{Stabilizer component groups.}
A weight of an exterior power is a sum of at most $d_\rho$ original
weights.  A difference of two such weights has entries at most
$2d_\rho M$ in the reciprocity coordinates.  A lattice generated by
such differences has rank $\ell\le r$.  The torsion order of its
quotient in the ambient free integral lattice is a gcd of maximal
minors, so it is at most
\[
       \ell!(2d_\rho M)^\ell\le(2d_\rho rM)^r=C_\Phi.
\]
This is the bound specified in \cite[\S2.3, before Table~1]{Jaswal}.
Hence $c_U=1/(2C_\Phi)$ is admissible.

The remaining formulas are now the ones in that table: one may take
$a=|\operatorname{GL}_{d_\rho}(\mathbf F_3)|$ and
$r_U=2\operatorname{rank}G=2r$, using its local-unit bound; they give
$B$, $B'$, $f$ and $\mu_{\rm unif}$ as displayed.
The chosen adjoint model is reductive at every $p>3$.  Explicitly,
after quadratic \emph{\'etale} splitting its adjoint representation is
that of $\operatorname{PGL}_m$ on $\mathfrak{sl}_m$.
Since $m$ is invertible over $\mathbf Z[1/6]$,
$\operatorname{Mat}_m=\mathcal O\,1\oplus\mathfrak{sl}_m$;
the closed algebra-automorphism group of $\operatorname{Mat}_m$
therefore embeds as a closed subgroup of
$\operatorname{GL}(\mathfrak{sl}_m)$.  Descent gives the same statement
for the unitary form.  Thus $C_0=4$ is compatible with both quadratic
fields, including the Eisenstein ramification at $3$.

Finally, at a relative-Weyl special point the adjoint Mumford--Tate
group is a maximal torus and $m_0=1$.  Its adjoint weights are its
$m(m-1)$ roots and zero weights.  The zero weights may be omitted from
the character list in Jaswal's discriminant argument: they contribute
neither valuations nor field generators.  The roots span the rational
character space, and their Hodge coefficients are $-1,0,1$.
Hence the proof of \cite[Proposition~3.3]{Jaswal} for this point uses
$t_\rho=1$ and $s_\rho=m(m-1)$.  The subsequent orbit and Hecke arguments
continue to concern this same special point.  This gives
\eqref{eq:ea-point-mu}; it does not replace the uniform higher-dimensional
statement by an unjustified sharper exponent.
\end{proof}

For convenient substitution the main numerical parameters are
\[
\begin{array}{c|r|r}
 m&4&8\\\hline
 \dim S=n^2&4&16\\
 {[F:\mathbf Q]}&2&2\\
 h_f&64&32768\\
 d_\rho&15&63\\
 r&3&7\\
 N&48&80640\\
 M&8&128\\
 C&13824&463613493490447548416\\
 NC&663552&37385792115069690304266240\\
 f&256&4096\\
 r_U&6&14\\
 C_0&4&4\\
 \mu^{-1}_{\rm unif}&4331520&839008101335040\\
 \mu^{-1}_{\rm point}&433152&5826445148160
\end{array}
\]
For $m=4$ one further has
\[
 A=(663552)!\,8\,(40320)^4,\quad
 B_\rho=18345885697,\quad C_\Phi=373248000.
\]
For $m=8$,
\[
 \begin{gathered}
 A=(37385792115069690304266240)!\,128\,(128!)^4,\\
 C_\Phi=233749065589177456392773613876412416.
 \end{gathered}
\]
and the short formula in \eqref{eq:ea-integer-parameters} specifies
$B_\rho$ exactly.  The Baily--Borel constants already used in the
forcing argument are $e^*(4)=15$, $e^*(16)=150$ and
$c_B=e^*(\dim S)^{2\dim S}$.

These choices evaluate the fixed group, level, field and finite-map
constants.  An application to a different prescribed level must account
for its finite level-change map.  An application to cycle detection
still requires the stated total-degree bound $H(D)$ and its projective
to Baily--Borel comparison.  Neither is supplied by the numerical
normalization alone.

\section{An explicit degree bound for the principal unitary family}
\label{sec:explicit-theta-chow-bound}

We use the Baily--Borel extension of the canonical bundle as the
polarization on the base.  Cycle degrees on an abelian variety are
$\theta^n[Z]$, without division by $n!$.  These two normalizations
are essential in the degree comparison.

We prove Theorem~\ref{thm:explicit-theta-chow-bound}.  In this
section $N_\theta$ denotes the dimension of theta projective space;
it is unrelated to the splitting-field degree $N=2g!$.


The integers are very large.  Their purpose is to eliminate unknown
equation degrees and unknown degrees of a map between compactifications.
The proof uses equations for arbitrary effective cycles, including
cycles with repeated components.  It therefore does not require a
claim that cubic integrability equations alone define the full Chow
variety.

\subsection{A relative Chow bound with specified degrees}

We first prove a general projective calculation.  Let
$X\subset\mathbf P^{N_\theta}_x\times\mathbf P^{N_\theta}_y$ be defined, over an
algebraically closed field of characteristic zero, by equations
$f_a(x,y)$ of bidegree $(2,2)$.  Additional equations on the parameter
$y$ will be imposed later.  Fix $r\geq2$ and $d\geq1$, and set
\begin{align*}
 c_d&=\binom{N_\theta+d}{d}^{r+1},&
 \ell_d&=(N_\theta+1)^d,& h_d&=(r+1)d,\\
 t_d&=\ell_d(h_d-1)+1,&
 b_d&=\binom{\ell_d+t_d-1}{t_d}.
\end{align*}

\begin{lemma}\label{lem:explicit-relative-chow}
In the projective space $\mathbf P^{c_d-1}$ of polynomials
$F(u_0,\ldots,u_r)$ of degree $d$ in each of the $r+1$ blocks of
$N_\theta+1$ variables, the pairs $(F,y)$ for which $F$ is the Chow form of
an effective $r$-cycle of degree $d$ supported on $X_y$ form a closed
set defined by bihomogeneous equations of bidegrees at most
$(b_d,2b_d)$.
\end{lemma}

\begin{proof}
We spell out a characterization of Chow forms used in
\cite[\S5.2, pp.~16--17]{GuerraPirola2009}.  For each choice of
$u_1,\ldots,u_r$, require the existence of points
$x_1,\ldots,x_d\in\mathbf P^{N_\theta}$ with the following properties:
\begin{enumerate}
\item the polynomials $F(u_0,u_1,\ldots,u_r)$ and
$\prod_{j=1}^{d}u_0(x_j)$ are proportional as polynomials in $u_0$;
\item $f_a(x_j,y)=0$ for every $a,j$;
\item $u_i(x_j)=0$ for all $1\leq i\leq r$ and $1\leq j\leq d$;
\item $F(v_0,\ldots,v_r)=0$ whenever all the $v_i$ pass through
one of the points $x_j$.
\end{enumerate}
Here proportionality is expressed by the $2$ by $2$ minors of the
two coefficient vectors.  This permits the first polynomial to be
zero.  The second is nonzero because each $x_j$ is a projective
point.  Points may repeat.

For clarity, this characterization includes nonreduced effective
cycles.  For a Chow form, the conditions hold for general
$u_1,\ldots,u_r$ by taking the intersection cycle with those
hyperplanes.  They hold for all choices because the incidence
conditions are closed and the points $x_j$ range over a projective
space.  Conversely, define $V_F$ to be the set of points $x$ such
that $F$ vanishes on every $(r+1)$-tuple of hyperplanes through $x$.
This is closed.  It has dimension at most $r$: otherwise every
$(r+1)$-tuple of hyperplanes would meet $V_F$, forcing $F=0$.
The asserted factorization for general $u_1,\ldots,u_r$ shows that
$V_F$ has an $r$-dimensional component.  Each such component has
its Chow form as an irreducible factor of $F$.  For general
$u_1,\ldots,u_r$, every zero of $F$ in the variable $u_0$ comes
from one of their intersection points with these components.
Consequently every factor of $F$ that depends on $u_0$ is one of
these Chow forms, with a positive multiplicity.  Their total degree
in $u_0$ is $d$, so their degrees already exhaust degree $d$ in
every block.  There is no remaining factor depending only on
$u_1,\ldots,u_r$.  The conditions $f_a(x_j,y)=0$ for general
linear sections imply that every component of the resulting cycle
is contained in $X_y$.  This proves the characterization.

We now bound its equations.  For an alternating matrix $A$, the
covector $x^tA$ passes through $x$, and every covector through a
nonzero $x$ has this form.  Expanding
\[
 F(x^tA_0,\ldots,x^tA_r)
\]
in the entries of the alternating matrices gives polynomials
$q_\alpha(F,x)$ that are linear in the coefficients of $F$ and
homogeneous of degree $h_d$ in $x$.  The fourth condition is
$q_\alpha(F,x_j)=0$ for all $\alpha,j$.

All four sets of incidence equations are homogeneous in each
$x_j$, of degree at most $h_d$.  Multiply each by all monomials
needed to give it multidegree $(h_d,\ldots,h_d)$ in the $d$
point blocks.  This does not change their zero set on
$(\mathbf P^{N_\theta})^d$.  Embed that product by Segre in
$\mathbf P^{\ell_d-1}$.  A polynomial of multidegree
$(h_d,\ldots,h_d)$ is the restriction of a degree-$h_d$
polynomial in Segre coordinates.  Adjoin the quadratic Segre
relations, multiplied by all monomials of degree $h_d-2$.
We have obtained a finite collection of homogeneous polynomials
$G_a$ of common degree $h_d\geq3$ in $\ell_d$ variables.  Their
coefficients are multihomogeneous in the remaining variables.
Each $G_a$ has degree at most one in $F$ and at most two in $y$.
The incidence projection over $(F,y,u_1,\ldots,u_r)$ is nonempty
exactly when the $G_a$ have a common projective zero.

Let $A=\mathbf C[z_1,\ldots,z_{\ell_d}]$ and
$J=(G_a)$.  The degree-$t_d$ Macaulay matrix is the matrix of
\[
 \bigoplus_a A_{t_d-h_d}\longrightarrow A_{t_d},\qquad
 (P_a)_a\longmapsto\sum_aP_aG_a.
\]
It has $b_d$ rows.  The common projective zero set is empty if and
only if this matrix has rank $b_d$.  One direction is immediate.
For the other, when there is no projective zero, choose
$\ell_d$ sufficiently general linear combinations of the $G_a$.
Successively avoiding the finitely many projective components
of the preceding intersection makes their common affine zero
set the origin.  They form a homogeneous system of parameters,
hence a regular sequence in the polynomial ring $A$.
The quotient by these combinations has Hilbert series
\[
 (1+z+\cdots+z^{h_d-1})^{\ell_d}.
\]
Consequently $A/J$ vanishes in all degrees at least
$t_d=\ell_d(h_d-1)+1$, as required.
This reasoning also covers specializations at which
some of the $G_a$ become zero.

It follows that nonemptiness is expressed by all maximal minors
of this matrix.  Each has degree at most $b_d$ in $F$ and at most
$2b_d$ in $y$.  Each column has a fixed multidegree in all the
parameter blocks, so every minor is multihomogeneous.  Requiring
the minors to vanish for every $u_1,\ldots,u_r$ amounts to
equating all their coefficients in these variables to zero.
This operation preserves the two stated degree bounds.  The
resulting equations are exactly the required relative Chow
equations.
\end{proof}

\begin{remark}
The common-degree condition makes the nonemptiness argument
particularly simple.  No effective Nullstellensatz estimate is
needed.  The much larger choice $t=h^\ell$ would also work by
the binomial theorem.
\end{remark}

\subsection{Theta coordinates and the integral unitary action}

Put $q=8$.  Mumford's moduli theorem gives an open subscheme
$M_q$ of a closed scheme
\[
 \overline M_q\subset\mathbf P^{q^g-1}
\]
defined by Riemann's quartic relations and the linear symmetry
relations.  It also gives a closed family
\[
 \overline A_q\subset\mathbf P^{q^g-1}\times\overline M_q
\]
defined over the base by equations of bidegree $(2,2)$; its
restriction to $M_q$ is the universal abelian scheme embedded
by the totally symmetric normalized line bundle of class
$q\theta$.  These are the theorem on p.~83 and the Further
Theorem on p.~84 of \cite{Mumford1967II}.  Only the existence of
the open subscheme is used.  The quartic scheme is not being
identified with a normal compactification.

We make the level convention precise.  A totally symmetric line
bundle in a given algebraic equivalence class is unique
\cite[p.~307]{Mumford1966I}.  Thus a principally polarized
abelian variety has a unique normalized totally symmetric bundle
of class $q\theta$; no choice of a symmetric representative of
$\theta$ adds a marking.  Conversely, a bundle of type
$(q,\ldots,q)$ has a Riemann form divisible by $q$, with
principal quotient.  A symmetric theta structure on this bundle
induces a symplectic marking on $A[q]$, and is induced by a
symplectic marking on $A[2q]$,
\cite[\S2, Remarks~2--4, pp.~318--320]{Mumford1966I}.

The stabilizer of this intermediate marking is exactly
\[
 \Gamma_g(q,2q)=
 \left\{\begin{pmatrix}A&B\\C&D\end{pmatrix}\in\Gamma_g(q):
 \operatorname{diag}(AB^t)\equiv
 \operatorname{diag}(CD^t)\equiv0\pmod{2q}\right\}.
\]
Here is a way to check that there is no unspecified covering
index in this identification.  Modulo $2q$, an element of
$\Gamma_g(q)$ is $I+q\left(\begin{smallmatrix}a&b\\c&-a^t
\end{smallmatrix}\right)$, with $b,c$ symmetric modulo $2$.
The diagonal blocks induce the identity on the theta structure
at level $q$.  Mumford's computation immediately after
Remark~3 on p.~319 says that the two off-diagonal blocks change
it by the sign character determined by
$\operatorname{diag}(b)$ and $\operatorname{diag}(c)$.
These signs vanish exactly when the two displayed congruences
hold, and all sign characters occur.  This gives the analytic
moduli space $\Gamma_g(q,2q)\backslash\mathcal H_g$ for the
same symmetric-theta-structure functor.  See also the explicit
identification in \cite[\S8A, p.~199]{LeFourn2019}; we do not
use the larger auxiliary coordinate system in that paper.
The group notation and its full symplectic action are also
given in \cite[\S1, p.~3]{FreitagSalvatiManni2016}.
For even $q$ the displayed congruence subgroup is normal in
$\operatorname{Sp}_{2g}(\mathbf Z)$ and contains
$\Gamma_g(2q)$.

One can see directly that these are the theta coordinates being
used here.  In the standard analytic trivialization their basis is
\[
 \vartheta_a(z,\tau)=
 \sum_{m\in\mathbf Z^g}
 \exp\!\left(\pi i q(m+a/q)^t\tau(m+a/q)
       +2\pi i q(m+a/q)^tz\right),
 \qquad a\in(\mathbf Z/q\mathbf Z)^g.
\]
The translations by the two standard $q$-torsion subgroups
act on this basis by permutations and diagonal characters,
respectively.  It is therefore the Schr\"odinger basis singled
out by the marking and by
\cite[Proposition~3 and Theorem~2, pp.~295--298]{Mumford1966I}.
At $z=0$ these are exactly the $f_{a,q}$ of
\cite[Introduction, p.~1]{FreitagSalvatiManni2016}.
The symplectic generators act projectively on these $q^g$
coordinates by permutations, by
$\vartheta_a(0,\tau+B)=\exp(\pi i a^tBa/q)\vartheta_a(0,\tau)$
for integral symmetric $B$, and by the finite Fourier matrix
$(\exp(-2\pi i a^tb/q))_{a,b}$ for
$\tau\mapsto-\tau^{-1}$.  Poisson summation gives the last
formula up to a common scalar.  In particular the resulting
projective matrices are algebraic and independent of $\tau$.

Let $J\in\operatorname{Sp}_{2g}(\mathbf Z)$ be multiplication by
$i$ or by $\omega$ on the prescribed lattice, and let $T_J$ be
its projective theta matrix.  The fixed locus of $T_J$ is defined
by the quadratic equations
\[
 (T_Jy)_a y_b-(T_Jy)_b y_a=0.
\]
Write $W_J$ for its intersection with $\overline M_q$.
The chosen connected unitary component at theta level, denoted
$S_\theta$, is open in $W_J$.

Here is a direct verification of this last assertion, including
the possible issue of extra automorphisms at special points.
On the Siegel upper half space, the locus fixed by $J$ is the
unitary domain for the given Hermitian lattice; for signature
$(n,n)$ it is connected and has dimension $n^2$.
Let $\Gamma=\Gamma_g(q,2q)$.  It is torsion-free and is
normalized by $J$.  The inverse image of the fixed locus of $J$
on $\Gamma\backslash\mathcal H_g$ is the union of the fixed
loci of elements in the coset $\Gamma J$.  These loci are
disjoint: if two elements in that coset fix the same point,
their quotient belongs to $\Gamma$ and to its finite point
stabilizer, so they coincide.  The union is locally finite by
proper discontinuity.  The image of the connected locus fixed
by $J$ is consequently an open and closed component of the
fixed locus on $M_q$; analytic and algebraic connected components
agree for this complex algebraic fixed locus.  Its arithmetic group is
$\Gamma\cap Z_{\operatorname{Sp}}(J)$, by the same stabilizer
argument.  This is exactly $S_\theta$.  Since $M_q$ is open in
$\overline M_q$, it follows that $S_\theta$ is open in $W_J$.
This argument permits additional automorphisms of an individual
abelian variety; it uses only the torsion-freeness of the level
group.

There is no extra determinant index in identifying this group
with the prescribed determinant-balanced unitary family.
An integral isometry has $K$-determinant in $\mathcal O_K^\times$,
and these units are roots of unity.  Every nontrivial such root
$\zeta$ satisfies $N_{K/\mathbf Q}(\zeta-1)\leq4$.
Thus congruence to one modulo $3$, or modulo $8$, forces that
determinant to be one.  At both levels the arithmetic group
therefore lies in the special unitary group, as required.

\subsection{The polarization comparison in the required direction}

Let $\lambda$ be the determinant of the Hodge bundle and let
$H$ be the hyperplane class of the theta-null map on $M_q$.
Then
\[
 H=\tfrac12\lambda\quad\hbox{in }\operatorname{Pic}(M_q)\otimes\mathbf Q.
\]
One can obtain this identity either from the determinant-bundle
formula or directly from Grothendieck--Riemann--Roch.
For the direct calculation, write $\pi:A\to M_q$ for the
universal family and $L$ for its normalized symmetric marking
line bundle.  Symmetry and normalization give
$[m]^*L\simeq L^{m^2}$.  Comparison with the degree $m^{2g}$
of $[m]$ shows that
$\pi_*c_1(L)^{g+1}=0$ rationally.  The relative tangent bundle
is the pullback of the dual Hodge bundle.  Since
$\operatorname{rank}\pi_*L=q^g$, Riemann--Roch gives
\[
 c_1(\pi_*L)=-\frac{q^g}{2}\lambda.
\]
The theta marking gives $\pi_*L=V\otimes M$ for a line bundle
$M$, by uniqueness of the theta-group representation.
Evaluation at the identity, where the normalized $L$ is trivial,
identifies $M$ with the pullback of $\mathcal O(-1)$ under the
theta-null map; this is the projective construction in
\cite[pp.~80--82]{Mumford1967II}.
Thus its first Chern class is $-q^gH$, proving the identity.

The theta-null map extends everywhere on the Siegel
Baily--Borel compactification.  Its normalization description
and the fact that it need not give a normal image are recalled
in \cite[Introduction, pp.~1--2]{FreitagSalvatiManni2016};
see also the normalization correction in
\cite[\S11, p.~228]{Mumford1967III}.  The absence of a base
locus can be seen on each rational boundary component: at a
standard boundary component, theta constants with zero indices
in the degenerating directions restrict to the theta constants
of the lower-dimensional abelian variety.  Those have no common
zero at its origin.  The finite Fourier transformation law
handles the other boundary components.  The map is finite,
since its hyperplane bundle is a positive rational multiple of
the ample Hodge bundle.  The Shimura map from the unitary base
extends to Baily--Borel compactifications
\cite[\S3.3]{DawJavanpeykarKuhne2020}, so the same theta
map is available on $S_\theta^{\rm BB}$.

We compute the unitary canonical class, rather than importing
the Siegel canonical weight.  On the unitary base the Hodge
bundle splits into its $K$-eigensubbundles
$E_\sigma\oplus E_{\bar\sigma}$, both of rank $n$.
The polarization identifies $E_{\bar\sigma}$ with
$(V_\sigma/E_\sigma)^\vee$ for the rank-$2n$ flat bundle
$V_\sigma$.  The cotangent bundle of the unitary domain is
$E_\sigma\otimes(V_\sigma/E_\sigma)^\vee$.  Therefore
\begin{align*}
 \det\Omega^1_{S_\theta}
 &= (\det E_\sigma)^n\otimes\det(V_\sigma/E_\sigma)^{-n}\\
 &=\bigl(\det(E_\sigma\oplus E_{\bar\sigma})\bigr)^n
 =\lambda^{\otimes n}.
\end{align*}
The Baily--Borel boundary of a signature-$(n,n)$ unitary
component has codimension at least $2n-1\geq3$.  On a normal
variety, restriction of line bundles to the complement of a
codimension-at-least-two set is injective.  Apply this after
clearing the denominators in the rational Picard identities.
For the Baily--Borel canonical bundle $\mathcal L_{\rm BB}$
and the extended theta map $\vartheta$ we obtain
\[
 c_1(\mathcal L_{\rm BB})
   =2n\,\vartheta^*c_1(\mathcal O(1))
   =g\,\vartheta^*c_1(\mathcal O(1)).
\]
This also proves directly that the restricted theta map is
finite.  The moduli theorem identifies $M_q$ itself with its
open projective theta image, with the same marking checked
above.  Thus this finite map is an embedding on $S_\theta$;
in particular it has generic degree one on every subvariety
of the open base.  For every
irreducible closed $Z\subset S_\theta$ of dimension $j$,
the projection formula gives the exact equality
\begin{equation}\label{eq:unitary-theta-degree}
 \deg_{\rm BB}Z=g^j\deg\overline{\vartheta(Z)}.
\end{equation}
The closure on the right is in theta projective space.  No
normality of that closure is required.

\subsection{Bounding the detector image and changing the level}

For a cycle of theta degree at most $D$, its degree in the
fiber embedding by $8\theta$ is at most $k=8^nD$.
For each $1\leq d\leq k$, apply
Lemma~\ref{lem:explicit-relative-chow} with $r=n$ to
$\overline A_q$.  Add the quartic equations of
$\overline M_q$ and the quadratic equations of $W_J$.
The resulting closed set in
$\mathbf P^{c_d-1}\times\mathbf P^{N_\theta}$ is defined by equations
of bidegrees at most $(2b_d,2b_d)$.  Its restriction over
$S_\theta$ is the required relative Chow space.  Because
$S_\theta$ is open in $W_J$, the closures of its irreducible
components are irreducible components of this closed set.
Multiply each equation by all monomials needed to reach
bidegree $(2b_d,2b_d)$; this preserves its zero set on the
product.  After the Segre embedding, these are equations of
degree $2b_d$; adjoin the quadratic Segre relations.
The total-degree B\'ezout bound gives total degree at most
\[
 (2b_d)^{(N_\theta+1)c_d-1}.
\]
All quantities increase with $d$, so the sum for
$1\leq d\leq k$ is at most $k e^R$.

The degree of the image of a projective parameter component
under projection to the base is at most its Segre degree.
Indeed, intersect its image with general base hyperplanes,
and then intersect the resulting fibers with general Segre
hyperplanes; the latter contribute a positive integer over
every point of the former.  The nonnegativity of mixed
intersection degrees gives the stated inequality.  A
detector is present on a union of connected components of
the relative cycle space, by the local constancy of the
projected cycle class.  Properness and the same argument
therefore bound the total theta degree of its image by
$k e^R$.  Repeated images are counted only once.  Applying
\eqref{eq:unitary-theta-degree}, and using $0\leq j\leq s$,
gives the theta-level bound $g^s k e^R$.  For the definition
using irreducible detectors, the image is unchanged: an
effective sum with nonzero projected class has an
irreducible component with nonzero projected class and no
larger degree.

Finally take the common cover of $S_\theta$ and the
principal-$3$ component $S$.  Its degree over $S_\theta$
is at most
\[
 |\operatorname{Sp}_{2g}(\mathbf F_3)|=I_g,
\]
because the corresponding index is the size of an image of
the theta-level arithmetic group modulo $3$.  Principal
level $48$ is contained in the intersection, so both maps
are finite level changes.  On their Baily--Borel
compactifications the canonical bundles pull back to the
canonical bundle of the common cover
\cite[\S3.2]{DawJavanpeykarKuhne2020}.  The cycle detector
condition is unchanged by level structure.  Over the open
base these finite maps are etale.  Thus for any irreducible
component of a detector locus the sum of the degrees of
its inverse-image components is its degree times the
covering degree, by the projection formula.  Mapping the
common cover to $S$ can only decrease total degree, since
each image component is counted once and each mapping
degree is a positive integer.  Multiplying the preceding
bound by $I_g$ proves
Theorem~\ref{thm:explicit-theta-chow-bound}.

The calculation is geometric.  The auxiliary theta level
and its projective matrices may be defined over a larger
number field, but the detector locus on the original
principal-$3$ family retains its original field of
definition.  No field-extension factor is added to the
geometric degree bound.


The same bound holds when total degree is defined as
$\sum_Z\max(1,\lceil\deg_{\rm BB}Z\rceil)$ over geometric
irreducible components.  Indeed, on the common theta and
principal-$3$ cover the degree of the closure of each component
meeting the open base is an integer: it is a positive integral
covering degree times $g^{\dim Z}$ times a theta projective degree.
If such a component $\widetilde Z$ maps finitely onto a component
$Z$ on the principal-$3$ base, the projection formula gives
$\deg_{\rm BB}\widetilde Z=m\deg_{\rm BB}Z$ for an integer
$m\geq1$.  Positivity and integrality therefore imply
\[
 \max(1,\lceil\deg_{\rm BB}Z\rceil)
 \leq\deg_{\rm BB}\widetilde Z.
\]
Choose one covering component for each distinct image component;
the components so chosen are distinct.  The preceding bound on
the total degree upstairs thus bounds this rounded total degree
downstairs as well.


\begin{corollary}\label{cor:theta-degree-whole-base}
With the same notation, the degree of the entire connected base
is at most
\[
 B_g=I_g\,g^s\,4^{N_\theta-s}.
\]
For $g=4$ this is
$131569513308979200\cdot4^{4095}$.
This bounds a detector locus only when that locus is the entire
base.  It does not bound an arbitrary proper sublocus.
\end{corollary}
\begin{proof}
The closure of the connected theta-level base is an
$s$-dimensional component of a scheme in $\mathbf P^{N_\theta}$ cut out
by equations of degree at most four.  The weighted B\'ezout argument in
Proposition~\ref{prop:bezout-degree}, before its weights are dropped,
gives $\deg(T)4^s\leq4^{N_\theta}$ for each such component $T$.
Hence its degree is at most $4^{N_\theta-s}$.  Apply the canonical factor $g^s$
and the same level-change bound $I_g$.
\end{proof}

\section{A degree certificate from equations}
\label{sec:degree-certificate}

This section identifies the finite projective calculation required to
evaluate the degree of a detector locus.  It also distinguishes the
exceptional components of a parameter scheme from fibres without detectors.

Choose a projective compactification $\overline S_0\subset\mathbf P^N$
of the fine base.  A projective embedding of its degree-$D$ cycle
parameter space over $S_0$ gives a graph compactification
\[
 \Gamma_D\ \subset\ \mathbf P^M\times\overline S_0.
\]
Take its reduced structure.  Its projection over $S_0$ is the given
cycle parameter morphism.  Write $a,b$ for the hyperplane classes from
$\mathbf P^N,\mathbf P^M$, respectively, and put
\[
 E_D=\sum_{T}\int_T(a+b)^{\dim T},
\]
where $T$ runs over all geometric irreducible components of $\Gamma_D$.
Thus $E_D$ is its total degree in the Segre embedding.  It can be
computed from homogeneous equations by irreducible decomposition and
Hilbert polynomials.  No decision whether a component detects a Weil
class is needed in computing this upper bound.

\begin{proposition}[Image degree and the Baily--Borel comparison]
\label{prop:degree-certificate}
Suppose that a projective realization of the Baily--Borel map on
$S_0$ is given by homogeneous forms of common degree $r$ in the
chosen coordinates, without a common zero on $S_0$.  Suppose its
hyperplane bundle is $\mathcal L_{\rm BB}^{\otimes e}$, with
$e\geq1$, and the map is finite on $S_0$.  If $s=\dim S_0$, then
\[
 \operatorname{tdeg}_{\rm BB}\mathcal D_D
 \leq\max\{1,(r/e)^s\}E_D.
\]
An integer upper bound is therefore
\[
 H(D)=\left\lceil\max\{1,(r/e)^s\}E_D\right\rceil.
\]
Here $\mathcal L_{\rm BB}$ is an integral ample bundle in the
normalization used in the discriminant theorem.  For a rational
normalization, first multiply by its fixed denominator.
\end{proposition}

\begin{proof}
Let $T$ be a component of $\Gamma_D$, let $Z$ be its image in
$\overline S_0$, and put $k=\dim T$ and $j=\dim Z$.
Intersecting $Z$ with $j$ general hyperplanes gives $\deg Z$
points in its dense image.  Intersecting their inverse images with
$k-j$ general Segre hyperplanes gives a positive number of points
above each.  Consequently
\[
 \deg Z\leq\int_T a^j(a+b)^{k-j}
            \leq\int_T(a+b)^k.
\]
The last inequality follows from the nonnegativity of the mixed
degrees of the two globally generated hyperplane bundles.
Every irreducible component of the detector image is the image of
one of these $T$.  Counting an image only once and summing proves
$\operatorname{tdeg}_{\mathbf P^N}\mathcal D_D\leq E_D$.

For an irreducible $j$-dimensional $Z\subset S_0$, pull back $j$
general hyperplanes from the Baily--Borel projective realization.
Their equations have degree $r$.  B\'ezout bounds the number of
isolated intersection points in $Z$ by $r^j\deg\overline Z$.
Any base points of the forms on the boundary can only increase the
total intersection bound; they are not counted as points in $Z$.
The finite map has positive degree on each image, which gives
\[
 e^j\deg_{\rm BB}Z\leq r^j\deg\overline Z.
\]
Since $0\leq j\leq s$, summing and using the first bound proves
the assertion.
\end{proof}

The direction of this comparison matters.  A bound for a theta
embedding in terms of Baily--Borel degree does not give the reverse
bound needed here.  The displayed equations of the map, or an
equivalent comparison of ample bundles, supply that reverse bound.

\begin{proposition}[A bound using only equation degrees]
\label{prop:bezout-degree}
Suppose $\Gamma_D\subset\mathbf P^M\times\mathbf P^N$
is cut out set-theoretically by bihomogeneous equations of
bidegrees $(a_i,b_i)$, where $a_i,b_i\leq e_D$.
Put
\[
 R_D=(M+1)(N+1)-1,\qquad e_D^+=\max(2,e_D).
\]
Then
\[
 E_D\leq(e_D^+)^{R_D},\qquad
 H(D)=\left\lceil\max\{1,(r/e)^s\}
                       (e_D^+)^{R_D}\right\rceil
\]
is an admissible degree certificate in
Proposition~\ref{prop:degree-certificate}.
\end{proposition}
\begin{proof}
Multiply an equation of bidegree $(a_i,b_i)$ by every
monomial of bidegree $(e_D-a_i,e_D-b_i)$.
Their common zero set on the product is the original zero
set: at each point some such monomial is nonzero.
Every resulting form of bidegree $(e_D,e_D)$ is the restriction
of a degree-$e_D$ form in the Segre coordinates.  Adjoin the
quadratic Segre equations.  Thus the Segre image is defined
set-theoretically in $\mathbf P^{R_D}$ by forms of degrees
at most $e_D^+$.

For completeness, a projective set $V\subset\mathbf P^R$
defined by equations of degrees at most $E$ has total degree
at most $E^R$.  After multiplying by monomials, use equations
of common degree $E$.  Start with $\mathbf P^R$ and weight an
irreducible component $T$ by $\deg(T)E^{\dim T}$.
Retain a component once it is contained in $V$.  Otherwise
intersect it with a general linear combination of the defining
equations which vanishes identically on none of the finitely
many components still to be cut.  The field is infinite, so
such a combination exists.  B\'ezout bounds the sum of the
weights of the resulting components by the preceding weight.
After at most $R+1$ steps every retained component is in $V$
and every other component has disappeared; a zero-dimensional
component outside $V$ disappears at the next step.
Every irreducible component of $V$ is among the retained
components, since the intersections preserve the containment
of $V$ at each step.  The initial weight was $E^R$, and each
final weight is at least its degree.  Duplicate retained
components can only enlarge this upper estimate.  This proves
the total-degree bound and then both displayed assertions.
\end{proof}

This bound avoids irreducible decomposition.  It still requires
the dimensions and equations of the chosen cycle parameter
space and the projective Baily--Borel map.  Those data, rather
than only the dimension of an abelian variety, determine its
numerical value.

\begin{definition}[Parameter exceptional set]\label{def:exceptional}
Let $f_D:\mathscr H_D^W\to\mathcal S$ be the proper parameter
space of detecting cycles on an integral arithmetic base.  Its
component exceptional set is
\[
 E_D^{\rm comp}=\bigcup_{T:\,\overline{f_D(T)}\ne\mathcal S}f_D(T),
\]
where $T$ runs over the irreducible components of
$\mathscr H_D^W$.  This is a finite union of closed subsets.  The
failure locus is the different set
$\mathcal S\setminus f_D(\mathscr H_D^W)$.
\end{definition}

A point outside $E_D^{\rm comp}$ which has a detector belongs to
the image of a dominating component.  Properness then forces a
detector on every fibre.  There is also a simpler certificate when
the generic detector locus is already the whole base.

\begin{proposition}[Known uniform bounds and arithmetic closure]
\label{prop:known-fourfolds}
Assume the degree-$48$ bound for the principal even Gaussian
fourfold family, or the degree-$36$ bound for the principal
Eisenstein fourfold family.  On any integral arithmetic model on
which the polarized family, labelled action, and proper cycle
parameter space exist, the image of the corresponding detector
parameter space is the entire base.  The failure locus is empty.
This assertion does not say that $E_D^{\rm comp}$ is empty, or
that the full parameter morphism is flat everywhere.
\end{proposition}

\begin{proof}
The assumed complex bound makes the proper image contain the
geometric generic fibre.  Its image is closed and therefore
contains the closure of the generic point of the integral model,
which is the entire model.  The universal projected cycle class
is defined on the open part where a chosen auxiliary prime
$\ell$ is invertible.  Apply this construction on $S[1/\ell]$.
Two distinct auxiliary primes give open parts covering $S$, so the
same argument covers every positive-characteristic fibre.  An irreducible detecting component of a parameter cycle
has no larger degree.
\end{proof}

The fourfold bounds are proved in \cite[Theorem~10.1,
pp.~19--22]{Bounded}, following van Geemen and
Schoen.  The Gaussian degree $48$ is the even-lattice bound; the
general odd-lattice bound there is $192$.  The cyclic Jacobian
calculation in that article gives $48$ on a particular odd-lattice
locus.  These statements already make the failure locus empty in
their respective ranges, without a discriminant computation.

For the eightfold family, Proposition~\ref{prop:degree-certificate}
provides an alternative certificate from a chosen projective model.
The individual integers $r,e,E_D$ of that alternative have not been
computed here. Theorem~\ref{thm:explicit-theta-chow-bound} already
gives the uniform explicit integer $H_g(D)$ without those choices.

\section{An explicit certificate for the exceptional primes}

The purpose of this section is to replace an unspecified finite
set by a finite algebraic calculation for a fixed fibre.

\begin{lemma}[A certificate for flatness over a number ring]
\label{lem:certificate}
Let $B=R[z_1,\ldots,z_m]/I$ be an explicitly presented algebra
of finite type over $R=\mathcal O_L[1/b]$. Its torsion ideal
$T=\ker(B\to B\otimes_RL)$ is killed by a nonzero element
of $\mathcal O_L$. A flat localization can be obtained by clearing denominators
in a monic Gr\"obner basis over $L$, together with certificates
that this basis generates the original ideal and that all its
$S$-polynomials reduce to zero. A power of that denominator
annihilates $T$.
\end{lemma}
\begin{proof}
The ideal $T$ is finitely generated because $B$ is noetherian.
Each generator is killed by some nonzero element of $R$; a
product kills all of $T$. Clearing the powers of $b$ makes that
product an algebraic integer. After its inversion $B$ is
torsion-free over a Dedekind domain, and is therefore flat.

For the computational form, compute a monic Gr\"obner basis
$g_1,\ldots,g_s$ for $I_L$. Keep polynomial identities expressing
each $g_i$ in the original generators, each original generator
in the $g_i$, and the zero remainders of Buchberger reductions.
There are finitely many coefficients in these identities. Invert
one algebraic integer containing all their denominators, including
the coefficients of the $g_i$. The identities now hold over
the localized ring, the leading coefficients remain one, and
the usual division algorithm shows that the standard monomials
span the quotient. If a linear combination of standard monomials
were zero, extension to $L$ and uniqueness of the normal form
would make every coefficient zero. They are thus a free basis,
possibly infinite. This proves flatness and the claimed effective
construction of a flatness denominator. To obtain an annihilator
of $T$ itself, compute generators of the saturation ideal and
raise that denominator to a power that kills all their images.
Such a power exists by finite generation. The first power need
not annihilate the torsion. No decision about algebraicity of a cohomology
class is used in computing this certificate.
\end{proof}

\begin{proposition}[The component exceptional set at a fixed fibre]
\label{prop:exact-exceptional}
For the fixed arithmetic parameter space of
Theorem~\ref{thm:arithmetic}, write
$\mathscr H_D=\bigcup_j\Spec B_j$ and set
\[
 \mathcal E_D^{\rm vert}
 =\left\{\mathfrak p\ne0:\ \mathfrak p=
           \mathfrak q\cap R\text{ for some }
           \mathfrak q\in\operatorname{Min}(B_j),
           \ \mathfrak q\cap R\ne0\right\}.
\]
This is precisely the set of places supporting a vertical
irreducible component of $\mathscr H_D$.
Together with the bad-model places it is the component
exceptional set used in the fixed-fibre lifting criterion.
It is finite and computable by primary decomposition of the
chart ideals and contraction to the number ring.
It is contained in the support of every valid flatness
denominator in Lemma~\ref{lem:certificate}.
\end{proposition}
\begin{proof}
Irreducible components of an affine noetherian scheme are the
closed sets of its minimal prime ideals.  A component is
vertical over a Dedekind base exactly when the contraction of
its generic prime is a nonzero prime of the base.  Every global
component meets at least one chart, and its intersection with
that chart is a component there.  This proves the formula
and finiteness.  Primary decomposition and elimination over
the explicitly presented number ring compute these
contractions.  After inversion of a flatness denominator,
flatness over the integral base excludes a vertical component:
a nonzero base element is then a nonzerodivisor on every chart.
This proves the containment.
\end{proof}

The containing set obtained from a flatness denominator can
be larger: embedded torsion also obstructs flatness.
Likewise, existence of a detector in every fibre does not
exclude other vertical components of the full parameter
space.  Thus the known fourfold bounds make the detector
failure locus empty, but by themselves give no assertion that
$\mathcal E_D^{\rm vert}$ is empty.

\begin{proof}[Proof of Theorem~\ref{thm:arithmetic}]
Apply Lemma~\ref{lem:certificate} on the finite affine cover.
After inverting $a_D$, the parameter scheme is flat over the
number ring. Pass to the strict henselization at the chosen
geometric place. Flatness of finite presentation makes every
nonempty open neighbourhood of a parameter point meet the
generic fibre; equivalently that point has a characteristic-zero
generization. Smooth proper base change identifies the cohomology
of the fixed family on this strictly henselian trait. The Betti
class, transported to $\ell$-adic cohomology through the chosen
embedding and geometric place, becomes a constant section there.
Universal cycle classes commute with specialization, so
\eqref{eq:oneprime-eq} holds at the generization with the same
$M$ and the same degrees.

The connected component of the geometric generic parameter
scheme containing the generization is defined over an algebraic
extension of $L$ and has an algebraic point. Since its ambient
variety is fixed, its cycle classes are constant along that
component. Choose such a point over $\overline L$. Comparison
between rational Betti and $\ell$-adic cohomology gives the
required rational equality, because scalar extension from
$\Q$ to $\Q_\ell$ is injective. The detector version uses
the same argument with the algebraic projector of
Lemma~\ref{lem:projector}; nonzero projected classes stay nonzero.
\end{proof}

\begin{corollary}[A degree-escape bound]\label{cor:escape}
For the fixed variety and class of Theorem~\ref{thm:arithmetic},
let $b_0$ be the minimum of
$\max(\deg Z^+,\deg Z^-)$ over representations with an arbitrary
positive integral multiplier in characteristic zero.  Define $b_v$
by the same minimum at a fixed geometric good reduction, allowing
the multiplier and every admissible auxiliary prime to vary. If $b_0=+\infty$,
then
\begin{equation}
 b_v>D\quad\text{whenever}\quad
 \Nm(v)>\bigl|\Nm_{L/\Q}(ba_D)\bigr|.
 \label{eq:escape}
\end{equation}
For certificates in each degree, set
\[
 A_D=\max_{1\leq e\leq D}
  \log\max\{2,|\Nm_{L/\Q}(ba_e)|\}.
\]
Every integer $D$ with $A_D<\log\Nm(v)$ is a strict lower
bound for $b_v$. This gives a computable degree-escape function
from the explicit parameter equations and their certificates.
\end{corollary}
\begin{proof}
If $v$ divides a nonzero integral principal ideal $(ba_D)$,
its norm is at most the norm of that ideal. Thus the displayed
norm inequality excludes every prime in the certificate. If
$b_v\leq D$, Theorem~\ref{thm:arithmetic} would produce a
characteristic-zero representative, contradicting $b_0=+\infty$.
The statement using $A_D$ follows by its definition.
\end{proof}

The bound depends on coefficients as well as geometric degrees.
This dependence cannot be removed by a general argument about
parameter schemes: the projective morphism
$\Spec\mathbf F_P\to\Spec\Z$ has a single geometric component
of degree one, for an arbitrarily large prime $P$. Its defining
coefficient and its exceptional prime both grow with $P$.
Equation~\eqref{eq:escape} gives a stopping threshold for a fixed
degree test. It does not give a finite procedure for declaring
an arbitrary class nonalgebraic when all degrees are allowed.

\section{Relative flatness and the choice of a prime}

This section states exactly what must be certified in the
moduli problem, in addition to complete splitting of a prime.

\begin{proof}[Proof of Theorem~\ref{thm:flat}]
The nonzero locus of the projected universal cycle section is
open and closed in $\mathscr H_D$, by the same argument as
Lemma~\ref{lem:closed}. Denote it by $\mathscr H_D^W$.
It is proper over $\mathcal S$. A morphism of finite presentation
which is flat at a point is open on a neighbourhood in its flat
locus. Therefore the image of $\mathscr H_D^W$ contains a
nonempty open subset of $\mathcal S$. Its proper image is closed.
Because $\mathcal S$ is integral, that image is all of
$\mathcal S$. In every fibre a detecting effective cycle has
an irreducible detecting component of no larger degree.
\end{proof}

Flatness over the number ring does not imply the flatness in
this theorem. The closed immersion
$\{u=0\}\to\mathbf A^1_{\Z}$ is flat over $\Z$. Its pullback
along the section $u=P$ is $\Spec\mathbf F_P$ and has no
characteristic-zero point. This example also shows why a list
of vertical primes of a universal parameter space cannot be
used unchanged for all CM sections.

\begin{proposition}[An abelian splitting-prime bound with exclusions]
\label{prop:prime}
Assume Dirichlet GRH. Let $L_0/\Q$ be abelian and contained
in $\Q(\zeta_f)$, and let $B$ be a finite set of forbidden
rational primes. Put
\[
 q=\operatorname{lcm}\left(4,f,\prod_{p\in B}p\right).
\]
There is a prime $p\notin B$, splitting completely in $L_0$,
with
\begin{equation}
 p<2(q\log q)^2.
 \label{eq:prime}
\end{equation}
The least such prime can be found by finite enumeration up
to this bound. This applies to forbidden model primes, the
residue characteristic of the chosen auxiliary prime, and the
rational primes below the support of an arithmetic certificate.
\end{proposition}
\begin{proof}
Bach--Sorenson \cite[Theorem 5.3, pp.~1731--1732]{BachSorenson}
gives a prime $p\equiv1\pmod q$ with the displayed bound under
Dirichlet GRH. Such a prime does not divide $q$, so it avoids
$B$. It splits completely in $\Q(\zeta_q)$ and hence in its
subfield $L_0$.
\end{proof}

Alp\"oge's asserted theorem gives precisely the analytic
hypothesis used in Proposition~\ref{prop:prime}. It also gives
the half-plane $\operatorname{Re}s>1/2$ for the Dedekind zeta
function of $L_0$, by its factorization into Dirichlet
$L$-functions. Proposition~\ref{prop:no-abelian-weyl} prevents using
that deduction for a relative-Weyl witness in the families at
issue. A nonabelian field requires an applicable effective
Chebotarev estimate and the same exclusions; a complete-splitting
condition alone does not certify relative flatness.

\begin{lemma}[Splitting the Mumford--Tate torus suffices]
\label{lem:split-ordinary}
Let $A$ be a CM abelian variety over a number field and let $L$
be the splitting field of its Mumford--Tate torus in its
faithful representation on $H^1(A,\Q)$.  If a rational prime
$p$ splits completely in $L$ and $A$ has good reduction at
a place above $p$, that reduction is ordinary.
\end{lemma}
\begin{proof}
Over $L$ the torus representation is a sum of its weight
spaces.  Its commutant is therefore a product of matrix
algebras over $L$.  By the description of Hodge endomorphisms
of an abelian variety, this commutant is
$\operatorname{End}^0(A)\otimes_\Q L$.
Its center is split over $L$.  If
$A\sim\prod_i A_i^{m_i}$ with the $A_i$ simple, that center
is $\prod_i E_i$, where $E_i$ is the CM field of $A_i$.
Consequently every $E_i$ splits at $p$.

Extend the field of definition so that the endomorphisms and
the isogeny are defined.  Newton slopes are unchanged by this
extension or by an isogeny.  For a simple factor, the
Shimura--Taniyama slope formula counts the members of its CM
type above a prime of $E_i$ and divides by the size of the
corresponding decomposition orbit.  Complete splitting makes
each orbit a singleton, so its slope is zero or one.
See \cite[Theorem~2.1, p.~3]{ConradST} for this formula.
Thus every factor, and hence $A$, is ordinary.  The argument
uses the rational CM action and so is unaffected by a
nonmaximal integral endomorphism order.
\end{proof}

Ordinarity gives no deformation of a specified cycle in the
Hilbert scheme; the arithmetic or relative-flatness certificate
is still needed.


The fixed-fibre arithmetic argument uses all good places outside
its certificate; it has no splitting restriction. The
specialization criterion for a family uses any prescribed
infinite set of good places. If a construction is only made
at split primes, its bound has only been established there.
For $\Q(i)$ the inert primes have $p\equiv3\pmod4$; for
$\Q(\omega)$ they have $p\equiv2\pmod3$. No detector-degree
bound at those primes follows merely by changing the congruence
in Proposition~\ref{prop:prime}.

\subsection{Combining the two certificates}

The arithmetic certificate can replace the unknown flatness over
the moduli base when the discriminant criterion is available.

\begin{corollary}[One prime at a relative-Weyl point]
\label{cor:combined-oneprime}
Use the fine unitary family and normalized adjoint level of
Corollary~\ref{cor:polynomial-detector}.  Let $x$ be a
relative-Weyl CM point, and choose a number field $L$ over which
its polarized fibre and labelled action have a smooth projective
model over $\mathcal O_L[1/b]$.  Fix $D$, compute a flatness
certificate $a_{x,D}$ for its degree-$D$ cycle parameter space,
and suppose
\[
 |\operatorname{disc}L_x|>
 \mathcal B_{S,F,\delta}(H(D))^{1/\mu}.
\]
Assume the analytic and effective Hecke inputs of that corollary.
If one good place $v\nmid ba_{x,D}$ admits a geometric detector
of degree at most $D$, then every complex fibre of the family
has a detector of degree at most $D$.
\end{corollary}

\begin{proof}
Theorem~\ref{thm:arithmetic} lifts the detector to the fixed
characteristic-zero fibre at $x$.  The splitting-field identity
in Lemma~\ref{lem:balanced-splitting-field} identifies $L_x$ with
the adjoint field in Corollary~\ref{cor:polynomial-detector}.
That corollary makes the closed detector locus the entire base.
\end{proof}

This corollary has no ordinarity hypothesis.  Complete splitting
is a useful way to choose a prime at which an additional geometric
construction might be available, but the lifting step uses the
flatness certificate for the fixed fibre.  It remains necessary to
produce the cycle with the stated degree at that prime.

\begin{proposition}[Nonabelian splitting with exclusions]
\label{prop:hu-exclusions}
Let $L_0/\mathbf Q$ be finite Galois, and let $B$ be a finite set
of forbidden rational primes.  Put
\[
 m=\operatorname{lcm}\left(4,\prod_{p\in B}p\right),\qquad
 L'=L_0\mathbf Q(\zeta_m),\qquad
 \mathcal D(L')=\log(e|d_{L'}|)+[L':\mathbf Q].
\]
Assume Hu's effective Chebotarev theorem
\cite[Corollary~9.6, pp.~56--57]{Hu}.  A prime splitting
completely in $L_0$ and avoiding $B$ exists with
\[
 p\leq C_{\rm Hu}\mathcal D(L')^8,
\]
where $C_{\rm Hu}$ is the effective absolute constant furnished
by that theorem.  Alternatively, enumerating unramified splitting
primes terminates.  The cited corollary does not print a numerical
value of $C_{\rm Hu}$.
\end{proposition}

\begin{proof}
Apply Hu's least-prime consequence to the identity conjugacy
class in $\operatorname{Gal}(L'/\mathbf Q)$.  Its conclusion
provides an unramified rational prime splitting completely in
$L'$.  Every odd prime in $B$ ramifies in
$\mathbf Q(\zeta_m)$, and $2$ ramifies because $4\mid m$.
Ramification persists in $L'$, so the prime is outside $B$.
It splits in the subfield $L_0$.  The bound and its effective
constant are those of the cited corollary.  The same conclusion
without an evaluated cutoff follows from enumerating primes and
factoring a defining polynomial away from its discriminant.
\end{proof}

The constant in Proposition~\ref{prop:hu-exclusions} must be
evaluated before its bound is a numerical least-prime cutoff.
The abelian result \eqref{eq:prime} already has its numerical
coefficient.  Neither result bounds the degree of a cycle.


\subsection{Relative-Weyl CM points and fixed ordinary reductions}
\label{subsec:weyl-ordinary-rigidity}

Throughout this subsection, discriminants are discriminants of number
fields, rather than discriminants of possibly nonmaximal orders.
We make the relative-Weyl hypothesis precise because the natural
permutation action is essential to the argument.
Let $A$ be a polarized abelian variety of balanced Weil type for an
imaginary quadratic field $K$, of dimension $m=2n\geq4$, and suppose
that $A$ is of CM type.  Write $T=\operatorname{MT}(A)$.
We call $A$ \emph{full relative-Weyl} if the image $T^{\mathrm{ad}}$
in the adjoint unitary group is a maximal torus and the Galois action
on its root system is the full natural action
\[
 \operatorname{Aut}(A_{m-1})=S_m\times\{\pm1\}.
\]
This is a condition on the action on the unitary root datum, not just
on the abstract isomorphism class of a finite Galois group.

\begin{lemma}[Signed permutation discriminants]
\label{lem:weyl-signed-discriminant}
Let $L/\mathbb Q$ be Galois with group
$G=S_m\times\{\pm1\}$, $m\geq4$, and let
$E=L^{S_{m-1}\times\{1\}}$, where $S_{m-1}$ is a natural point
stabilizer.  Put $N=2\cdot m!$.
Then
\begin{equation}
 \label{eq:weyl-discriminant-divisibility}
 |d_L|\mid |d_E|^{N/4}.
\end{equation}
In particular $|d_L|\leq |d_E|^{N/4}$, with no restriction on wild
ramification.
\end{lemma}

\begin{proof}
Identify $G/(S_{m-1}\times\{1\})$ with
$X=\{1,\ldots,m\}\times\{\pm1\}$, on which
$(s,\epsilon)(i,\eta)=(s(i),\epsilon\eta)$.
Every nonidentity element moves at least four points of $X$:
an element with $\epsilon=-1$ fixes no point, while a nontrivial
permutation with $\epsilon=1$ moves at least two letters in each
of the two signed copies.  A transposition attains the minimum four.
For any subgroup $J\subset G$, Burnside's lemma therefore gives
\[
 \#(J\backslash X)
 =\frac{1}{|J|}\sum_{g\in J}|X^g|
 \leq 2m-4\bigl(1-|J|^{-1}\bigr).
\]
For the permutation representation $V=\mathbb C[X]$ and the regular
representation $R=\mathbb C[G]$, this says
\begin{equation}
 \label{eq:weyl-all-subgroups}
 \operatorname{codim} R^J
 =N\bigl(1-|J|^{-1}\bigr)
 \leq\frac N4\operatorname{codim} V^J
 \qquad(J\subset G).
\end{equation}

At a rational prime $p$, let $I_i$ be the lower ramification groups
at a chosen prime of $L$ above $p$.  The Artin conductor formula is
\[
 a_p(U)=\sum_{i\geq0}\frac{|I_i|}{|I_0|}
                      \operatorname{codim}U^{I_i}.
\]
All coefficients are nonnegative, so
\eqref{eq:weyl-all-subgroups} implies
$a_p(R)\leq (N/4)a_p(V)$, including at wildly ramified primes.
The conductor--discriminant identity gives
$a_p(R)=v_p(|d_L|)$ and $a_p(V)=v_p(|d_E|)$.
Since $N/4=m!/2$ is an integer, this proves
\eqref{eq:weyl-discriminant-divisibility} prime by prime.
For the conductor--discriminant identity and the equivalent positive
expansion in upper ramification groups, see
\cite[\S3C, equations (3-5)--(3-7), pp.~614--615]{JonesRoberts14WO};
the comparison is also the present signed-action specialization of
\cite[Theorem~7.3, p.~634]{JonesRoberts14WO}.
\end{proof}

\begin{lemma}[The CM field at a full relative-Weyl point]
\label{lem:weyl-cm-field}
Let $A$ be a full relative-Weyl balanced Weil CM abelian variety of
complex dimension $m\geq4$, as above. Then $A$ is absolutely simple,
and $E=\operatorname{End}^0(A_{\overline{\mathbb Q}})$ is a CM field
of degree $2m$. If $L$ is the splitting field of its Mumford--Tate
torus, then
\[
 \operatorname{Gal}(L/\mathbb Q)=S_m\times C_2,\qquad
 E=L^{S_{m-1}\times\{1\}},\qquad L=E^{\rm nor}.
\]
The only subfields of $E$ are $\mathbb Q,E^+,K,E$; in particular,
$K$ is its only proper CM subfield. No choice of reduction is needed.
\end{lemma}
\begin{proof}
Write $\nu$ for the rational similitude character.  The two
$K$-eigenspaces in $H^1(A,\mathbb Q)\otimes\overline{\mathbb Q}$
have weights
\begin{equation}
 \label{eq:weyl-standard-weights}
 \lambda_i=\tfrac12\nu+u_i,\qquad
 \overline\lambda_i=\tfrac12\nu-u_i,
 \qquad \sum_{i=1}^{m}u_i=0.
\end{equation}
Indeed, balanced Weil type gives the determinant relation
$\sum_i\lambda_i=n\nu$, and the roots are the differences
$u_i-u_j$.  The relative-Weyl action is
$(s,\epsilon)u_i=\epsilon u_{s(i)}$.
The roots together with $\nu$ span the rational character space of
$T$: the kernel of the adjoint map has connected part contained in
the scalar torus, since the balanced determinant relation removes
the connected norm-one central factor.  Consequently the splitting
fields of $T$ and $T^{\mathrm{ad}}$ agree.  The Galois group of $L$
is therefore $G=S_m\times C_2$ and $[L:\mathbb Q]=N$.
Complex conjugation acts as $u_i\mapsto-u_i$, since the
polarizable CM torus is compact modulo its weight torus; hence it
is the central element $(1,-1)$.

For $m\geq3$, the $2m$ weights in
\eqref{eq:weyl-standard-weights} are distinct, with the transitive
faithful signed permutation action.  Thus the commutant of $T$ on
$H^1(A,\mathbb Q)$ is an etale algebra of degree $2m$, whose
primitive idempotents are permuted transitively by Galois.  It is
therefore a field $E$.  Rational Hodge endomorphisms of $H^1$ are
the rational endomorphisms of $A$, so
$E=\operatorname{End}^0(A_{\overline{\mathbb Q}})$.
The absence of nontrivial idempotents proves absolute simplicity.
The stabilizer of a weight is $H=S_{m-1}\times\{1\}$, giving
\begin{equation}
 \label{eq:weyl-field-stabilizer}
 E=L^H.
\end{equation}
Its normal core is trivial, so $L$ is the normal closure of $E$.
Complex conjugation is central and nontrivial on $E$, showing that
$E$ is CM, with $[E^+:\mathbb Q]=m$.  The given quadratic field
$K$ is fixed by $S_m\times\{1\}$, and $E=E^+K$.

To list the intermediate fields, let $J$ contain $H$ and project
$J$ to $S_m$.  The point stabilizer $S_{m-1}$ is maximal, so that
projection is either $S_{m-1}$ or $S_m$.  In the first case,
$J=H$ or $S_{m-1}\times C_2$.  In the second case, either $J=G$,
or $J$ is the graph of a homomorphism $S_m\to C_2$ whose restriction
to $S_{m-1}$ is trivial.  The sign homomorphism is nontrivial on
$S_{m-1}$ for $m\geq4$, so the only latter possibility is
$J=S_m\times\{1\}$.  Galois correspondence gives exactly the four
fields in the statement.

\end{proof}

\begin{theorem}[Ordinary simplicity and the splitting-field bound]
\label{thm:weyl-ordinary-rigidity}
Let $A$ be a full relative-Weyl balanced Weil CM abelian variety as
above, defined over a number field.  Let $L$ be the splitting field
of its Mumford--Tate torus and put $N=2\cdot m!$.
Then the following statements hold.
\begin{enumerate}
\item
$A$ is absolutely simple and
$E=\operatorname{End}^0(A_{\overline{\mathbb Q}})$ is a CM field of
degree $2m$.  Its normal closure is $L$,
$\operatorname{Gal}(L/\mathbb Q)=S_m\times C_2$, and the complete
list of subfields of $E$ is
\[
 \mathbb Q,\qquad E^+,\qquad K,\qquad E.
\]
In particular, $K$ is its only proper CM subfield.
\item
Every ordinary good geometric reduction $A_0$ is absolutely simple,
and specialization identifies
\[
 \operatorname{End}^0(A_{0,\overline{\mathbb F}_p})=E.
\]
If $A_0$ is defined over $\mathbb F_q$, its actual $q$-Frobenius
$\pi$ generates $E$ over $\mathbb Q$.
\item
For such an ordinary reduction over $\mathbb F_q$,
\begin{equation}
 \label{eq:weyl-ordinary-q-bound}
 |d_L|\leq(2\sqrt q)^{m(2m-1)N/2},
 \qquad
 q\geq\frac14 |d_L|^{\,4/(m(2m-1)N)}.
\end{equation}
For $m=4$ this gives $q\geq\tfrac14|d_L|^{1/336}$, and for
$m=8$ it gives $q\geq\tfrac14|d_L|^{1/2419200}$.
\end{enumerate}
There is no maximal-order hypothesis and no hypothesis that $p$ be
unramified in $E$ or $L$.
\end{theorem}

\begin{proof}
The first assertion is Lemma~\ref{lem:weyl-cm-field}.

\emph{Specialization and the balanced norm relation.}
Fix an ordinary good reduction $A_0/\mathbb F_q$ and a prime
$\ell\ne p$.  Good reduction gives an injection
$E\hookrightarrow\operatorname{End}^0(A_{0,\overline{\mathbb F}_p})$.
Smooth proper base change identifies the rational Tate module with
a free module of rank one over $E\otimes\mathbb Q_\ell$.
Its commutant in the geometric rational endomorphism algebra is
therefore exactly $E$: it contains $E$, and its faithful action on
the Tate module has dimension at most $2m$.

Choose a finite extension $\mathbb F_Q/\mathbb F_q$, $Q=q^r$,
over which all geometric endomorphisms of $A_0$ are defined.
Enlarge it if necessary as follows.  The balanced Weil plane
\[
 W_K=\bigwedge_K^m H^1(A,\mathbb Q)
 \subset H^{2n}(A,\mathbb Q)
\]
has Hodge type $(n,n)$.  Deligne's theorem makes its elements
absolute Hodge classes, and their Galois action becomes trivial
after a finite extension
\cite[Proposition~2.9(b), pp.~18--19, and Main Theorem~2.11,
p.~19]{DeligneAH82WO}.  After specialization and the Tate twist,
the $Q$-Frobenius therefore fixes $W_K(n)$ pointwise.
Let $\pi_r$ be that Frobenius endomorphism.  It commutes with $E$,
so $\pi_r\in E$, and its action on the $K$-determinant line is
\begin{equation}
 \label{eq:weyl-frobenius-norm}
 \operatorname{Nm}_{E/K}(\pi_r)=Q^n.
\end{equation}
For a smaller field the quotient can be a root of unity; the
finite extension removes it.  This step uses the absolute-Hodge
property of the Weil classes, not their algebraicity.

\emph{Ordinarity and geometric endomorphisms.}
The polarization gives $\pi_r\overline\pi_r=Q$, where the bar is
complex conjugation in $E$.  If $\pi_r\in E^+$, then
$\pi_r^2=Q$; every normalized $p$-adic valuation is $1/2$.
If $\pi_r\in K$, equation~\eqref{eq:weyl-frobenius-norm} gives
$\pi_r^m=Q^n$, again forcing every normalized valuation to be
$n/m=1/2$.  Both conclusions contradict ordinarity.  Since the
subfield list has no other proper possibility,
\[
 \mathbb Q(\pi_r)=E.
\]
Tate's endomorphism theorem now identifies
\[
 \operatorname{End}^0(A_0/\mathbb F_Q)\otimes\mathbb Q_\ell
 =\operatorname{End}_{\mathbb Q_\ell[\pi_r]}(V_\ell A_0)
 =E\otimes\mathbb Q_\ell;
\]
see \cite[\S1, Main Theorem, p.~134, and Lemma~1, p.~135]{Tate66WO}.
All geometric endomorphisms were defined over $\mathbb F_Q$, so
the geometric rational endomorphism algebra equals $E$.
Its lack of nontrivial idempotents again proves absolute
simplicity.  Finally the original $q$-Frobenius $\pi$ lies in this
field and $\pi^r=\pi_r$ generates $E$.  Thus $\pi$ itself generates
$E$.  This proves (2) with the original residue cardinality $q$.

\emph{Discriminants.}
The algebraic integer $\pi$ has $2m$ distinct conjugates
$\pi_1,\ldots,\pi_{2m}$, all of absolute value $\sqrt q$.
The field discriminant divides the discriminant of its power-basis
order, whence
\[
 |d_E|
 \leq\left|\operatorname{disc}(\operatorname{minpol}(\pi))\right|
 =\prod_{i<j}|\pi_i-\pi_j|^2
 \leq(2\sqrt q)^{2m(2m-1)}.
\]
Apply Lemma~\ref{lem:weyl-signed-discriminant} and raise the last
bound to $N/4$.  This gives
\eqref{eq:weyl-ordinary-q-bound}.  For $m=4$ the exponent on
$2\sqrt q$ is $672$; for $m=8$ it is $4838400$, giving the stated
reciprocal exponents on $|d_L|$.
\end{proof}

\begin{corollary}[A fixed ordinary seed fixes the splitting field]
\label{cor:weyl-fixed-ordinary-seed}
Fix a geometric ordinary abelian variety $B$ over a finite field.
Among full relative-Weyl balanced Weil CM lifts whose good
geometric reduction is isomorphic, or merely geometrically
isogenous, to $B$, the CM field $E$ and its normal closure $L$ are
fixed up to isomorphism.  In particular their splitting-field
discriminants cannot tend to infinity.
If $B$ is not absolutely simple, no such lift exists.
\end{corollary}

\begin{proof}
Geometric isogeny preserves the rational endomorphism algebra.
Theorem~\ref{thm:weyl-ordinary-rigidity} identifies this algebra
with the field $E$, and $L$ is its normal closure.
\end{proof}

\begin{remark}[Order discriminants and the detector problem]
\label{rem:weyl-order-scope}
For an order $\mathcal O\subset\mathcal O_E$ one has
\[
 |\operatorname{disc}(\mathcal O)|
 =|d_E|\,[\mathcal O_E:\mathcal O]^2.
\]
Specialization embeds the characteristic-zero endomorphism order
in the order of the reduction; the proof above does not bound the
index in this display.  A complexity involving order
discriminants therefore requires separate conductor control.
The corollary constrains a strategy that varies the splitting field
while retaining one ordinary seed.  It supplies no semiregular
object, no spreading family, and no nonzero Weil-class detector.
\end{remark}

\section{The orbit bound inherited from the Siegel moduli space}
\label{sec:orbit-transfer}

This section separates a bound in the discriminant of an endomorphism
order from a bound in the splitting field of a torus.
Let $F$ be a number field and let
$f:S\longrightarrow\mathcal A_g$ be an $F$-morphism.
For a point $x\in S(\overline{\mathbf Q})$ whose image is CM, put
\[
 D_x=\left|\operatorname{disc}
 Z\bigl(\operatorname{End}_{\overline{\mathbf Q}}A_{f(x)}\bigr)\right|.
\]

\begin{proposition}\label{prop:orbit-transfer}
Assume the Siegel-moduli orbit estimate
\[
 |\operatorname{Gal}(\overline{\mathbf Q}/\mathbf Q)y|
 \geq c_g D_y^{1/(4g)}
\]
for CM points $y\in\mathcal A_g(\overline{\mathbf Q})$.
Then every such $x$ satisfies
\[
 |\operatorname{Gal}(\overline{\mathbf Q}/F)x|
 \geq \frac{c_g}{[F:\mathbf Q]}D_x^{1/(4g)}.
\]
In particular this applies to the morphism to $\mathcal A_g$ attached
to a Shimura variety of Hodge type.  Neither injectivity nor finiteness
of $f$ is required.
\end{proposition}

\begin{proof}
Equivariance makes the image of the $F$-Galois orbit of $x$ the entire
$F$-Galois orbit of $f(x)$.  Thus the first orbit has at least as many
points as the second.  If $E=\mathbf Q(f(x))$, then
\[
 |\operatorname{Gal}(\overline{\mathbf Q}/F)f(x)|
 =[FE:F]=\frac{[FE:\mathbf Q]}{[F:\mathbf Q]}
 \geq\frac{[E:\mathbf Q]}{[F:\mathbf Q]}.
\]
Apply the assumed estimate to $f(x)$.
\end{proof}

The proof explains why no comparison of discriminants is a new input
to this restriction of the Siegel estimate.  Such a comparison is
needed only to express it in the different invariant used by
Jaswal's Proposition~3.3.  In particular the exponents
$1/(4g)$ and $1/(16tsn_T(n_T-1))$ cannot be compared numerically
without also identifying their discriminants and hypotheses.

\begin{example}\label{ex:conductor-torus}
Fix an imaginary quadratic field $K$.  For each positive integer $f$,
an elliptic curve with complex multiplication by the order
$\mathbf Z+f\mathcal O_K$ has centre discriminant
$f^2|d_K|$.  Its Mumford--Tate torus is split by $K$, independently of
$f$.  Hence the centre discriminant cannot be bounded in terms of
the torus splitting-field discriminant alone.  A level or conductor
invariant is indispensable in that direction.
\end{example}

\subsection{An explicit Siegel orbit constant from a zero-free half-plane}
\label{subsec:explicit-siegel-orbit}

We now prove the lower orbit estimate used above, with a constant
suited to the same analytic hypothesis as the discriminant criterion.
The ingredients are the averaged Colmez formula, the isogeny
estimates of Gaudron--R\'emond, and a direct zero-sum comparison.
This follows the height-to-orbit argument of Tsimerman
\cite[proof of Theorem~4.1 and Theorem~5.1]{Tsimerman}, using the
sharper discriminant estimate in \cite[Theorem~1.5(2)]{GR23}.
It requires no algebraicity theorem for CM Hodge classes.

\begin{lemma}\label{lem:halfplane-height}
Let $E$ be a CM field of degree $2r$, let
$D_E=|\operatorname{disc}E|$, and assume
$\operatorname{QRH}(\delta;E)$ with $0<\delta\leq1/2$.
For an abelian variety $A$ with complex multiplication by
$\mathcal O_E$, of any CM type, its stable Faltings height in
the normalization of \cite[\S2.3]{GR14} satisfies
\[
 h_F(A)\leq a_r\log D_E+b_{r,\delta},\qquad
 a_r=3\cdot2^{r-2},\quad
 b_{r,\delta}=\frac{a_r}{\delta}+4r(2^r-1).
\]
\end{lemma}
\begin{proof}
Put $t=1/\log D_E$ and $\sigma=1+t$; since $D_E\geq3$,
$1<\sigma<2$.  The logarithmic derivative of the completed
Dedekind zeta function, paired over its nontrivial zeros,
gives
\begin{align*}
 \sum_\rho\operatorname{Re}\frac1{\sigma-\rho}
 ={}&\tfrac12\log D_E+\frac1\sigma+\frac1{\sigma-1}
       +\frac{\zeta_E'}{\zeta_E}(\sigma)\\
    &+r\bigl(\psi(\sigma)-\log(2\pi)\bigr),
\end{align*}
where $\psi=\Gamma'/\Gamma$.  The logarithmic derivative of
$\zeta_E$ is negative.  Also
$\psi(\sigma)\leq\psi(2)=1-\gamma$, so
$1/\sigma+r(\psi(\sigma)-\log(2\pi))<0$.
Consequently the sum is less than $\tfrac32\log D_E$.
These are the same Hadamard-product identities used in
\cite[Lemma~4.2, pp.~12--14]{Wei}; here every nontrivial zero $\rho$ satisfies
$\operatorname{Re}\rho\leq1-\delta$, and no exceptional zero is
separated out.

For $\rho=\beta+i\gamma_\rho$ write $u=1-\beta\geq\delta$.
The nonnegative summands satisfy
\[
 \frac{\operatorname{Re}(1/(1-\rho))}
      {\operatorname{Re}(1/(\sigma-\rho))}
 =\frac{u}{u+t}
       \frac{(u+t)^2+\gamma_\rho^2}{u^2+\gamma_\rho^2}
 \leq\frac{u+t}{u}\leq1+\frac t\delta.
\]
Summing, first over a finite symmetric set of zeros and then
passing to its increasing limit, proves
\begin{equation}\label{eq:halfplane-zero-sum}
 \sum_\rho\operatorname{Re}\frac1{1-\rho}
 \leq\frac32\log D_E+\frac3{2\delta}.
\end{equation}

Let $E^+$ be the maximal totally real subfield and let
$L(s)=\zeta_E(s)/\zeta_{E^+}(s)$.
This is the entire quadratic Hecke $L$-function of $E/E^+$.
Its zeros form a submultiset of those of $\zeta_E$.
The functional equation of
\[
 \Lambda(s)=(D_E/D_{E^+})^{s/2}
     \bigl(\pi^{-(s+1)/2}\Gamma((s+1)/2)\bigr)^r L(s)
\]
and its Hadamard product give
\[
 -\frac{L'(0)}{L(0)}
 =\sum_{\Lambda(\rho)=0}\operatorname{Re}\frac1{1-\rho}
   +\frac12\log\frac{D_E}{D_{E^+}}
   +\frac r2\bigl(\psi(1/2)-\log\pi\bigr).
\]
The averaged Colmez formula, proved independently in
\cite[Theorem~1.1, p.~535]{YuanZhang} and
\cite[Theorem~A, p.~392]{AGHMP}, identifies
\[
 2^{-r}\sum_\Phi h_W(A_\Phi)
 =-\frac12\frac{L'(0)}{L(0)}
        -\frac14\log\frac{D_E}{D_{E^+}},
 \qquad h_W=h_F+\frac r2\log\pi.
\]
Yuan--Zhang use $h_W$ as defined above; the height in
Andreatta--Goren--Howard--Madapusi Pera is
$h_W-(r/2)\log(2\pi)$, and their displayed constant changes
accordingly.  The last gamma term is negative.  Thus the average is at most
$\tfrac34\log D_E+\tfrac3{4\delta}$.
The height lower bound \cite[Corollary~8.4]{GR14} implies
$h_W(A_\Phi)>-4r$ for every type.
Isolate the chosen summand among the $2^r$ types and use
$h_F\leq h_W$ to obtain the claimed constants.
\end{proof}


\begin{proposition}[Use the whole CM-type orbit]
\label{prop:cm-type-orbit-height}
Let $E$ be a CM field of degree $2r$ and absolute discriminant $D_E$,
and assume that $\zeta_E$ has no zeros in
$\operatorname{Re}(s)>1-\delta$, where $0<\delta\leq1/2$.
Let $\Phi$ be a CM type whose orbit under
$\operatorname{Gal}(\overline{\mathbb Q}/\mathbb Q)$ has cardinality
$t$.  If $A$ has CM type $\Phi$ and maximal $\mathcal O_E$ action,
then
\begin{equation}
 \label{eq:cm-type-orbit-height}
 h_F(A)\leq
 \frac{3\cdot2^{r-2}}{t}
       \bigl(\log D_E+\delta^{-1}\bigr)
 +4r\left(\frac{2^r}{t}-1\right).
\end{equation}
\end{proposition}

\begin{proof}
Use the normalization $h_W=h_F+(r/2)\log\pi$.
The averaged Colmez formula and the zero-sum calculation in the
proof of Lemma~\ref{lem:halfplane-height} give
\[
 2^{-r}\sum_{\Psi}h_W(E,\Psi)
 \leq\frac34\bigl(\log D_E+\delta^{-1}\bigr).
\]
Stable Faltings height is invariant under conjugating an abelian
variety by an automorphism of $\overline{\mathbb Q}$.  Conjugation
sends the CM type $\Phi$ to its Galois conjugate, so the $t$ types
in its orbit all have height $h_W(E,\Phi)$.  Each of the other
$2^r-t$ types has height at least $-4r$.  Therefore
\[
 t\,h_W(E,\Phi)
 \leq3\cdot2^{r-2}\bigl(\log D_E+\delta^{-1}\bigr)
       +4r(2^r-t).
\]
Divide by $t$ and use $h_F\leq h_W$.
\end{proof}

\begin{corollary}[Balanced full relative-Weyl types]
\label{cor:balanced-weyl-height}
Let $A$ be full relative-Weyl and of balanced Weil type, of dimension
$m=2n\geq4$, as in Theorem~\ref{thm:weyl-ordinary-rigidity}.
Let $E=\operatorname{End}^0(A)$, assume maximal $\mathcal O_E$
action, and impose the zero-free hypothesis of
Proposition~\ref{prop:cm-type-orbit-height} on $E$.
Then
\begin{equation}
 \label{eq:balanced-weyl-height}
 h_F(A)\leq
 \frac{3\cdot2^{m-2}}{\binom{m}{m/2}}
       \bigl(\log D_E+\delta^{-1}\bigr)
 +4m\left(\frac{2^m}{\binom{m}{m/2}}-1\right).
\end{equation}
In particular, the right sides for dimensions four and eight are
\[
 2\bigl(\log D_E+\delta^{-1}\bigr)+\frac{80}{3},
 \qquad
 \frac{96}{35}\bigl(\log D_E+\delta^{-1}\bigr)
       +\frac{2976}{35},
\]
respectively.
\end{corollary}

\begin{proof}
The $2m$ embeddings of $E$ are the signed letters
$(i,\epsilon)$, $1\leq i\leq m$, $\epsilon\in\{\pm1\}$.
A CM type chooses one sign for each letter.  Its signature with
respect to $K$ counts the chosen positive and negative signs.
Balanced signature therefore means choosing exactly $m/2$ positive
signs.  The natural $S_m$ action is transitive on these
$\binom{m}{m/2}$ choices, and the central sign change preserves this
set.  Thus the type orbit has exactly
$t=\binom{m}{m/2}$ elements.  Apply
Proposition~\ref{prop:cm-type-orbit-height} with $r=m$.
\end{proof}

\begin{remark}
Complex conjugation sends every CM type to a distinct type of the
same height.  Consequently $t\geq2$ even without the relative-Weyl
hypothesis; isolating that pair already halves the single-summand
coefficient in Lemma~\ref{lem:halfplane-height}.
For a nonmaximal endomorphism order, the displayed estimates apply
after saturation, with the corresponding isogeny-height cost added.
They do not remove the need to control that cost.
\end{remark}

For a positive integer $g$ define
\begin{align*}
 A_g&=3\cdot2^{g-2},&
 B_{g,\delta}&=A_g/\delta+4g(2^g-1),\\
 C_g&=241(eg)^{2g}g^5,& s_g&=2\cdot3^{4g^2},\\
 W_{g,\delta}&=g(A_g+\tfrac12)
       +\frac{gB_{g,\delta}+\tfrac32g+1}{\log3},&
 c_{g,\delta}&=\frac{e}{4g s_g C_g W_{g,\delta}}.
\end{align*}

\begin{theorem}[The explicit endomorphism-discriminant orbit bound]
\label{thm:explicit-siegel-orbit}
Let $y\in\mathcal A_g(\overline\Q)$ be a CM point, and put
$D_y=|\operatorname{disc}Z(\operatorname{End}_{\overline\Q}A_y)|$.
Assume $\operatorname{QRH}(\delta;E_i)$ for every CM field
occurring as a factor of the center of
$\operatorname{End}^0(A_y)$.  Then
\[
 |\Gal(\overline\Q/\Q)y|
 \geq\frac{D_y^{1/(2g)}}{s_gC_gW_{g,\delta}\log D_y}
 \geq c_{g,\delta}D_y^{1/(4g)}.
\]
Consequently Proposition~\ref{prop:orbit-transfer} and
Corollary~\ref{cor:centre-forcing} apply with
$c_g=c_{g,\delta}$.  Cureton--Thorner's theorem supplies
$\delta=1/32$, and Hu's supplies $\delta=1/8$, subject to
those respective analytic inputs.  If every $E_i$ is abelian,
the asserted Dirichlet GRH supplies $\delta=1/2$.
\end{theorem}

\begin{proof}
First let $A/K$ be a CM abelian variety with all geometric
endomorphisms defined over the number field $K$.  Write
\[
 A\sim\prod_i A_i^{m_i},\quad
 E_i=\operatorname{End}^0(A_i),\quad [E_i:\Q]=2r_i,
 \quad\sum_i m_ir_i=g,
\]
with simple pairwise nonisogenous factors.  Put
$E=\prod_i E_i$, $\mathcal O=\prod_i\mathcal O_{E_i}$,
$R=Z(\operatorname{End}A)$, and $I=[\mathcal O:R]$.
Then $D=|\operatorname{disc}R|=I^2\prod_iD_{E_i}\geq3$.

Choose a complex uniformization $A(\C)=V/\Lambda$.
Since $I\mathcal O\subset R$, the saturated lattice satisfies
$\Lambda\subset\mathcal O\Lambda\subset I^{-1}\Lambda$.
The quotient $\mathcal O\Lambda/\Lambda$ is a $K$-defined
subgroup of $A[I]$: for a basis $u_j$ of $\mathcal O$, it
is the sum of the images of $A[I]$ under the integral
$K$-endomorphisms $Iu_j$.
The quotient $A\to B=V/\mathcal O\Lambda$ is therefore a
$K$-isogeny of degree at most $I^{2g}$.
On each $E_i$-factor the lattice is a projective
$\mathcal O_{E_i}$-module.  Its decomposition into invertible
ideals decomposes the complex torus into simple CM factors
with maximal orders.  The projection idempotents are
endomorphisms of $B$, all defined over $K$, so this is also a
product decomposition over $K$.

Faltings height is additive on products and changes by at most
half the logarithm of the degree under an isogeny
\cite[\S2.3]{GR14}.  Lemma~\ref{lem:halfplane-height}, the
monotonicity of $a_r,b_{r,\delta}$ for $r\leq g$, and
$\log I\leq\tfrac12\log D$ therefore give
\begin{equation}\label{eq:order-height}
 h_F(A)\leq g(A_g+\tfrac12)\log D+gB_{g,\delta}.
\end{equation}

For a prime $\ell$ let
$S_\ell=\Z_\ell[\rho_\ell(\Gal(\overline K/K))]$.
Since every endomorphism is defined over $K$, the Galois
action commutes with $\operatorname{End}A$.  For a CM variety
its rational commutant on the Tate module is $E\otimes\Q_\ell$.
The integral form of Faltings's endomorphism theorem identifies
the Galois-equivariant endomorphisms of the Tate module with
$\operatorname{End}(A)\otimes\Z_\ell$; it follows as well from
the rational theorem by saturation.  Taking centers, or
intersecting with $E\otimes\Q_\ell$, gives the integral
commutant $R\otimes\Z_\ell$.  Thus
$S_\ell\subset R\otimes\Z_\ell$.
Faltings's endomorphism theorem and the double-centralizer
identity imply $S_\ell\otimes\Q_\ell=E\otimes\Q_\ell$;
equivalently one may use \cite[Proposition~13.7]{GR23}.
The intrinsic trace form on this commutative algebra gives
\[
 \operatorname{disc}S_\ell
  =\operatorname{disc}(R\otimes\Z_\ell)
          [R\otimes\Z_\ell:S_\ell]^2.
\]
Consequently $D$ divides $\prod_\ell\operatorname{disc}S_\ell$.

Apply Gaudron--R\'emond to the class generated by $A$.
In their Notations~1.1--1.2 the auxiliary dimension is
$n=\sum_i r_i\leq g$.
Their Theorem~1.5(2) and the inequality after Notation~1.2,
together with Theorem~1.8, give
\[
 D\leq d\leq\Upsilon^{2n},\qquad
 \Upsilon\leq C_g[K:\Q]
          \max\{1,\log[K:\Q],h_F(A)+\tfrac32g\}.
\]
These are statements about intrinsic discriminants; no
representation-trace multiplicity is inserted.
If $k=[K:\Q]\leq D$, equation~\eqref{eq:order-height}
makes the maximum at most $W_{g,\delta}\log D$.
Hence
\[
 k\geq\frac{D^{1/(2g)}}{C_gW_{g,\delta}\log D}.
\]
If $k>D$, the same inequality holds because its right side
is at most $D$.

Finally, start with the moduli point $y$.  A symplectic
level-three marking and a primitive cube root of unity give
a field of definition $K$ of degree at most
$s_g[\Q(y):\Q]$: there are fewer than $3^{4g^2}$ choices
of a torsion basis, and a level-three marked principally
polarized variety has no nontrivial automorphisms.
All geometric endomorphisms are defined over $K$.
Indeed their finite Galois action is trivial modulo three,
because $A[3]$ is rational, and the level-three congruence
subgroup of an integral general linear group has no torsion.
Applying the field-degree inequality to this $K$ proves
the first orbit bound.  The elementary inequality
\[
 \log D\leq\frac{4g}{e}D^{1/(4g)}
\]
proves the second.
\end{proof}

\subsection{A sharper point orbit bound and discriminant threshold}
\label{subsec:hybrid-point-forcing}

The displayed uniform constants below use the general
$c_{m,\delta}$ of Theorem~\ref{thm:explicit-siegel-orbit}.
The balanced CM-type refinement can improve that coefficient;
the leading exponent in $H$ is unchanged.


The two available orbit estimates control different quantities.
Jaswal's estimate contains the level invariant $\alpha_y$, whereas
Theorem~\ref{thm:explicit-siegel-orbit}, combined with
Lemma~\ref{lem:weyl-signed-discriminant}, gives a larger exponent of
the splitting discriminant.  We use both inequalities to bound the
prime needed for a Hecke intersection.  Their product is not used.
The argument also gives a sharper orbit upper bound than
Proposition~\ref{prop:orbit-upper} in the full relative-Weyl setting.

Throughout this subsection the abelian dimension is $m\in\{4,8\}$,
the unitary signature is $(m/2,m/2)$, and
\[
 N=2m!,\qquad s=m^2/4,\qquad q=2^{s-1},\qquad
 \nu=\frac1{mN},\qquad
 \mu=\frac1{16m(m-1)N(N-1)}.
\]
Use the fine component $S_0$, adjoint component $S$, finite
$F$-morphism $\pi_{\rm ad}:S_0\to S$ and fibre bound $h_f$ of
Proposition~\ref{prop:ea-fine-normalization}.
More generally the argument works after a fixed extension of their
field of definition; write $d=[F:\Q]$.
For the concrete normalization, $F=\Q(\zeta_3)$, $d=2$, and one may
take $h_f=64$ for $m=4$ and $h_f=32768$ for $m=8$.
All Baily--Borel degrees below are total geometric degrees.
Use the constants of Proposition~\ref{prop:ea-explicit-constants},
with the point exponent in \eqref{eq:ea-point-mu}; the definitions
of $\beta_0$ and $C_*$ use this same $\mu$.  Put
\[
 k=k_1d+3f,\qquad
 \mathcal R(H)=c_B^{q-1}H^q,\qquad h=k(q-1).
\]
Here $f=(\dim G+1)^2$ is the fixed datum constant.

\begin{lemma}[Point intersections over the original field]
\label{lem:hybrid-point-hecke}
Let $y\in S(\overline\Q)$ be a full relative-Weyl special point.
Let $Z\subsetneq S$ be closed and $F$-irreducible, contain $y$, and
have dimension $z$ and total degree $D\geq1$.
Suppose that $l>C_0$ splits completely in $L_y$, that
$K_{T_y,l}=K^m_{T_y,l}$, and that $l>1/c_U$.
Then
\begin{equation}\label{eq:hybrid-point-hecke}
 |\Gal(\overline\Q/F)y|
 \leq c_B^{2^z-1}D^{2^z}l^{k(2^z-1)}.
\end{equation}
For $z=0$ the same assertion holds without a condition on $l$.
\end{lemma}

\begin{proof}
When $z=0$, $Z$ contains the $F$-orbit of $y$, so its degree is
at least the orbit cardinality.  Suppose $z>0$.  By Lemma~\ref{lem:weyl},
every positive-dimensional geometric component through $y$ is
Hodge generic.  The same is true of all its $F$-conjugates and of
their finite-level lifts.

We extract the cutting operation in the proof of
\cite[Theorem~6.1, Steps~1--4]{Jaswal}, retaining its fields and
degree bounds.  At the prime $l$ it supplies a level
$I\subset K^{\rm ad}$, a finite map
$\pi:\operatorname{Sh}_I\to\operatorname{Sh}_{K^{\rm ad}}$, and a Hecke element $t_l\in T_y(\Q_l)$.
The cover has degree at most $l^f$ and the Hecke correspondence
$T_{t_l}$ at level $I$ has degree less than $l^{k_1d+f}$.
The construction uses a reciprocity power
$j\leq[FE_y:E_y]\leq d$, already included in the exponent
$k_1d+f$.  There is a lift $\widetilde y$ of $y$ and an element
$\sigma\in\Gal(\overline\Q/FE_y)$ such that
$\sigma\widetilde y\in T_{t_l}\widetilde y$.

Let $\widetilde Z$ be the union of the $F$-conjugates of a geometric
component of $\pi^{-1}(Z)$ containing $\widetilde y$.
It is $F$-irreducible, $\pi(\widetilde Z)=Z$, and
\[
 \deg\widetilde Z\leq l^fD,\qquad
 \deg T_{t_l}\widetilde Z<l^{k_1d+2f}D.
\]
Both $\widetilde Z$ and $T_{t_l}\widetilde Z$ are defined over $F$.
Since $\sigma$ fixes $F$, the preceding reciprocity relation implies
$\widetilde y\in\widetilde Z\cap T_{t_l}\widetilde Z$.
Their intersection therefore contains the entire $F$-orbit of
$\widetilde y$.

The stable alternative $\widetilde Z\subset
T_{t_l}\widetilde Z$ is impossible.  The argument of
\cite[Theorem~6.1, Step~3]{Jaswal}, using the unboundedness property
of the chosen Hecke element, then produces a positive-dimensional
special subvariety through $\widetilde y$ inside $\widetilde Z$.
Its image is a positive-dimensional special subvariety through $y$
inside $Z$.  By Lemma~\ref{lem:weyl} that image must be $S$,
contrary to $Z\subsetneq S$.  The monodromy hypothesis in this
stable-case argument is satisfied: each lifted geometric component
is Hodge generic and contains an $F$-conjugate of $\widetilde y$.

The geometric components of $\widetilde Z$ are transitively
permuted over $F$.  Since the Hecke image is defined over $F$,
failure of the stable alternative means that it contains none of
those components.  Every component of the intersection thus has
dimension less than $z$.  Choose one through $\widetilde y$, take
the union $\widetilde Z'$ of its $F$-conjugates, and put
$Z'=\pi(\widetilde Z')$.  This is a closed $F$-irreducible locus
through the original $y$, of smaller dimension.  The uniform
Baily--Borel B\'ezout inequality
\cite[Lemma~5.2]{Jaswal} and the projection formula give
\[
 \deg Z'\leq\deg\widetilde Z'
 \leq c_B l^{k_1d+3f}D^2.
\]
If $Z'$ has positive dimension, it is again Hodge generic by
Lemma~\ref{lem:weyl}.  We repeat this operation on $Z'$ at the
original level $K$, over the same field $F$, with the same point
$y$ and prime $l$.  Thus no new special closure, component field,
or accumulated change of base level occurs.

After at most $z$ cuts we obtain the zero-dimensional $F$-orbit of
$y$.  Iterating $D\mapsto c_Bl^kD^2$ gives
\[
 (c_Bl^k)^{2^z-1}D^{2^z},
\]
which proves \eqref{eq:hybrid-point-hecke}.  If the process stops
earlier, the same expression still bounds its degree because
$D,c_B,l^k\geq1$.
\end{proof}

In particular, if $y$ lies in a proper $F$-closed locus of total
degree at most $H$, take the $F$-conjugate union of one component
through $y$.  Its dimension is at most $s-1$.  Under
$\operatorname{QRH}(\delta;L_y)$, the prime supplied by
\cite[Proposition~6.6]{Jaswal} gives
\begin{equation}\label{eq:hybrid-geometric-orbit-upper}
 |\Gal(\overline\Q/F)y|\leq\mathcal R(H)x_y^h,
\end{equation}
where, writing $\beta=\log d_{L_y}$ and $\ell=\log\alpha_y$,
\begin{equation}\label{eq:hybrid-prime-cost}
 x_y=\max\left\{
 (64\kappa_\delta\beta)^{1/\delta},\,
 (\beta+1)^2,\,(48N)^2,\,4C_0,\,
 64N^2\bigl(\max(e^2/B',\ell)+2\bigr)^2
 \right\}.
\end{equation}
The zero-dimensional case is covered as well, since its orbit has
at most $H\leq\mathcal R(H)x_y^h$ points.
Compared with the displayed form of Jaswal's general criterion,
\eqref{eq:hybrid-geometric-orbit-upper} retains the exponent
$q-1$ of the Hecke cost and does not insert an extra $d^q$ into
the degree factor.  This follows from the point argument just
proved, not from substituting a new right-hand side into the
statement of his theorem.

For the next result use the explicit constant
$c=c_{m,\delta}$ of Theorem~\ref{thm:explicit-siegel-orbit} and set
\begin{equation}\label{eq:hybrid-orbit-coefficients}
 \mathsf a=\frac{C_*}{d},\qquad
 \mathsf b=\frac{c}{d h_f}.
\end{equation}
These numbers belong to $(0,1)$.  Define the following constants,
all independent of $H$ and of the point:
\begin{align}
 \lambda_\delta&=64\kappa_\delta+1,&
 M_N&=8N,\notag\\
 \Gamma_\delta&=1+48N+4C_0+8N(e^2/B'+2),\notag\\
 \mathsf K&=M_N+\frac{\lambda_\delta-M_N\mu}{\nu},
 \label{eq:hybrid-affine-constants}\\
 \mathsf C&=\Gamma_\delta+M_N\log(1/\mathsf a)
       +\frac{\lambda_\delta-M_N\mu}{\nu}
          \log(1/\mathsf b),\notag\\
 \tau&=\frac{(k_1d+3f)(q-1)}{\delta}.\notag
\end{align}
Here $\lambda_\delta-M_N\mu>0$, and hence
$\mathsf K,\mathsf C>0$.  Finally put
\begin{align}
 \mathcal V(H)&=2\log\mathcal R(H)+\frac{\mathsf C}{\mathsf K}
       +2\tau\log\frac{2\mathsf K\tau}{e},\notag\\
 \mathcal J(H)&=\mathcal R(H)
       \bigl(\mathsf K\mathcal V(H)+\mathsf C\bigr)^\tau,
 \label{eq:hybrid-orbit-threshold}\\
 \mathcal X(H)&=\frac{d h_f\mathcal J(H)}{c}
              =\frac{\mathcal J(H)}{\mathsf b}.
 \label{eq:hybrid-X}
\end{align}
For the stated data $2\mathsf K\tau/e>1$, so these are positive
closed expressions, with $\mathcal X(H)>1$ for $H\geq1$.

\begin{theorem}[The combined point orbit bound]
\label{thm:hybrid-point-orbit}
Let $\mathscr Y\subsetneq S_0$ be reduced, closed and defined over
$F$, with total Baily--Borel degree at most $H\geq1$.
Suppose that $x\in\mathscr Y(\overline\Q)$ is a full
relative-Weyl CM point and put $y=\pi_{\rm ad}(x)$.
Assume $\operatorname{QRH}(\delta;L_y)$ and the hypotheses of
Theorem~\ref{thm:explicit-siegel-orbit} for the CM centre field of
$A_x$.  Then
\begin{align}
 |\Gal(\overline\Q/F)y|&\leq\mathcal J(H),&
 |\Gal(\overline\Q/F)x|&\leq h_f\mathcal J(H),
 \label{eq:hybrid-actual-orbits}\\
 \alpha_y d_{L_y}^{\mu}&\leq\frac{\mathcal J(H)}{\mathsf a},&
 d_{L_y}&\leq\mathcal X(H)^{1/\nu}.
 \label{eq:hybrid-power-bound}
\end{align}
For the centre-order discriminant $D_x$ one also has
$D_x\leq\mathcal X(H)^{4m}$.
\end{theorem}

\begin{proof}
The image $\pi_{\rm ad}(\mathscr Y)$ is proper, closed, defined over $F$,
and of total Baily--Borel degree at most $H$, by finiteness and
the equality of Baily--Borel bundles under pullback.
Thus \eqref{eq:hybrid-geometric-orbit-upper} applies to $y$.
Write
\[
 O=|\Gal(\overline\Q/F)y|,\qquad u=\log O\geq0.
\]
The point form of \cite[Corollary~3.4 and Lemma~6.4]{Jaswal}
gives
\[
 O\geq\mathsf a\alpha_y d_{L_y}^{\mu}.
\]
The loss $d$ in $\mathsf a$ is sufficient because
$[FE_y:E_y]\leq d$ and the orbit over $FE_y$ is contained in the
orbit over $F$.

By Lemma~\ref{lem:weyl-cm-field}, the CM endomorphism field $E$
has the signed permutation action on its embeddings.  Its maximal-order discriminant
satisfies $d_E\leq D_x$.  Lemma~\ref{lem:weyl-signed-discriminant}
therefore gives
\[
 d_{L_y}\leq d_E^{N/4}\leq D_x^{N/4}.
\]
Here Lemma~\ref{lem:balanced-splitting-field} identifies the full
and adjoint torus splitting fields.  Theorem
\ref{thm:explicit-siegel-orbit} and Proposition
\ref{prop:orbit-transfer} imply
\[
 |\Gal(\overline\Q/F)x|
 \geq\frac c d D_x^{1/(4m)}
 \geq\frac c d d_{L_y}^{\nu}.
\]
The $F$-orbit of $x$ maps onto that of $y$, with at most $h_f$
points over each image, so $O\geq\mathsf b d_{L_y}^{\nu}$.
There is no passage to a reflex field in this second estimate.

The two lower bounds give, separately,
\[
 \ell+\mu\beta\leq u+\log(1/\mathsf a),\qquad
 \nu\beta\leq u+\log(1/\mathsf b).
\]
As $\lambda_\delta-M_N\mu>0$, they imply
\begin{align*}
 \lambda_\delta\beta+M_N\ell+\Gamma_\delta
 &\leq M_N\bigl(u+\log(1/\mathsf a)\bigr)
       +(\lambda_\delta-M_N\mu)\beta+\Gamma_\delta\\
 &\leq\mathsf K u+\mathsf C.
\end{align*}
Every term of \eqref{eq:hybrid-prime-cost} is at most
$(\lambda_\delta\beta+M_N\ell+\Gamma_\delta)^{1/\delta}$:
the quantity in parentheses bounds respectively
$64\kappa_\delta\beta$, $\beta+1$, $48N$, $4C_0$, and
$8N(\max(e^2/B',\ell)+2)$, and $1/\delta\geq2$.
Consequently
\[
 x_y\leq(\mathsf K u+\mathsf C)^{1/\delta}.
\]
Combining this with \eqref{eq:hybrid-geometric-orbit-upper} yields
\[
 u\leq\log\mathcal R(H)
       +\tau\log(\mathsf K u+\mathsf C).
\]
Set $w=u+\mathsf C/\mathsf K$.  Applying
$\log w\leq w/(2\tau)+\log(2\tau)-1$ to the right side gives
$u\leq\mathcal V(H)$.  Substitution into the preceding
inequality proves $O\leq\mathcal J(H)$.
The finite-fibre inequality gives the second bound in
\eqref{eq:hybrid-actual-orbits}; the two lower bounds give
\eqref{eq:hybrid-power-bound}.  Finally
$(c/d)D_x^{1/(4m)}\leq h_f\mathcal J(H)$ gives the stated bound
for $D_x$.
\end{proof}

\begin{corollary}[Using the logarithmic centre-discriminant bound]
\label{cor:hybrid-log-centre}
Under the hypotheses of Theorem~\ref{thm:hybrid-point-orbit},
\begin{align}
 D_x&\leq
       \bigl(e\mathcal X(H)\log\mathcal X(H)\bigr)^{2m},
 \label{eq:hybrid-centre-threshold}\\
 d_{L_x}=d_{L_y}&\leq
       \bigl(e\mathcal X(H)\log\mathcal X(H)\bigr)^{mN/2}.
 \label{eq:hybrid-splitting-threshold}
\end{align}
\end{corollary}

\begin{proof}
The first inequality of Theorem~\ref{thm:explicit-siegel-orbit},
transferred to $F$, and \eqref{eq:hybrid-actual-orbits} give
\[
 D_x^{1/(2m)}
 \leq d h_f\mathcal J(H)s_m C_m W_{m,\delta}\log D_x.
\]
The power bound just proved gives
$\log D_x\leq4m\log\mathcal X(H)$.  Since
$c=e/(4m s_m C_m W_{m,\delta})$, substitution yields
\[
 D_x^{1/(2m)}
 \leq e\mathcal X(H)\log\mathcal X(H).
\]
This proves \eqref{eq:hybrid-centre-threshold}.
Using $d_{L_y}\leq d_E^{N/4}\leq D_x^{N/4}$ proves
\eqref{eq:hybrid-splitting-threshold}.
\end{proof}

For fixed datum, field and $\delta$, the expressions have the sizes
\begin{align*}
 \mathcal J(H)&=H^{q+o(1)},\\
 \bigl(e\mathcal X(H)\log\mathcal X(H)\bigr)^{2m}
    &=H^{2mq+o(1)},\\
 \bigl(e\mathcal X(H)\log\mathcal X(H)\bigr)^{mN/2}
    &=H^{qmN/2+o(1)}.
\end{align*}
More precisely, $\mathcal J(H)$ is a fixed multiple of
$H^q$ times the $\tau$-th power of a positive affine function of
$\log H$.  Thus every exponent strictly larger than the displayed
principal exponent admits a closed polynomial upper bound, by
$\log H\leq H^\epsilon/(e\epsilon)$.
The principal splitting-discriminant exponents are
\[
\begin{array}{c|r|r|r}
 m&q/\mu&q/\nu&qmN/2\\ \hline
 4&3\,465\,216&1\,536&768\\
 8&190\,920\,954\,614\,906\,880&21\,139\,292\,160&
      10\,569\,646\,080
\end{array}
\]
The middle column uses only the power consequence of the centre-discriminant bound;
the last uses its first, logarithmic inequality.  The improvement
factors from $q/\mu$ to the last column are respectively
$4512$ and $18\,063\,136$.

\begin{corollary}[A stronger detector gate]
\label{cor:hybrid-detector}
Assume the total-degree certificate and detector hypotheses of
Corollary~\ref{cor:polynomial-detector}, together with the analytic
hypotheses of Theorem~\ref{thm:hybrid-point-orbit}.
In that corollary the splitting-discriminant threshold may be
replaced by
\begin{equation}\label{eq:hybrid-detector-gate}
 \min\left\{
 \mathcal B_{S,F,\delta}(H(D))^{1/\mu},\,
 \bigl(e\mathcal X(H(D))\log\mathcal X(H(D))\bigr)^{mN/2}
 \right\}.
\end{equation}
Alternatively a witness with centre-order discriminant greater
than \eqref{eq:hybrid-centre-threshold}, evaluated at $H(D)$,
forces the same conclusion.  The fine-orbit upper bound
$h_f\mathcal J(H(D))$ can replace $U(H(D))$ in
Corollary~\ref{cor:centre-forcing} under these hypotheses.
\end{corollary}

\begin{proof}
If the detector locus were proper, its total degree would be at
most $H(D)$.  Theorem~\ref{thm:fixed-field-bound} and Corollary
\ref{cor:hybrid-log-centre} would then impose both thresholds in
\eqref{eq:hybrid-detector-gate}, contradicting the strict witness
inequality.  The centre-discriminant and orbit formulations follow
from the same two results.
\end{proof}

The minimum retains the earlier threshold whenever it is smaller
at a given degree.  The new principal exponent is strictly smaller,
so the new bound is the smaller one for all sufficiently large
$H$.  This refinement still requires an actual detector, its
membership witness and a total base-degree certificate $H(D)$.
The Hecke argument uses the indicated proofs in Jaswal's
unrefereed paper and their named external inputs; Theorem~\ref{thm:explicit-siegel-orbit} and the
signed-permutation discriminant comparison are the separate inputs
already stated above.  The Gaussian eightfold detector remains open.

\subsection{An upper orbit bound inside a proper detector locus}
\label{subsec:orbit-upper}

An alternative upper bound, using only the original discriminant and
level inequalities, can be made explicit from the same
discriminant and level estimates.  Let $f:S_0\to S$ be the finite
$F$-map from the fine unitary component to its adjoint component,
and let $h_f$ bound the cardinality of every geometric fibre.
For the finite map of neat Shimura components one may take its
degree.  Put $r=\operatorname{rank}G$ and let $N\geq2$ bound the
degrees of the torus splitting fields.  Keep the constants
$B\leq1$, $B'=Bc_U$ and $C_0$ of
Section~\ref{sec:polynomial-threshold}.  All logarithms below
are natural, except for the displayed $\log_2$.

\begin{proposition}\label{prop:orbit-upper}
Set
\begin{align*}
 \lambda&=\log(1/B),& R&=r(N-1),\\
 p_{\rm orb}&=\max\left\{1+\lambda,
                \frac{r/2+\lambda/\log2}{\mu}\right\},\\
 C_{\rm orb}&=h_f N^{1+r(\log_2N+N)}2^{Nr}
                    B^{-(C_0+e^2/B')},\\
 U(H)&=C_{\rm orb}\mathcal B(H)^{p_{\rm orb}}
       \left(1+N\log2+\frac{\log\mathcal B(H)}{2\mu}\right)^R,
\end{align*}
where $\mathcal B(H)=\mathcal B_{S,F,\delta}(H)$.
Under the effective inputs of Theorem~\ref{thm:forcing}, a
relative-Weyl point $x\in S_0$ whose image lies in a proper
$F$-closed locus of total Baily--Borel degree at most $H$ satisfies
\[
 |\Gal(\overline\Q/F)x|\leq U(H).
\]
Consequently a detector at $x$ of degree at most $D$ together
with $|\Gal(\overline\Q/F)x|>U(H_D)$ forces
$\mathcal D_D=S_0$.
\end{proposition}

\begin{proof}
Write $y=f(x)$, let $T$ be its adjoint Mumford--Tate torus,
let $L$ split $T$, and put $\Delta=|\operatorname{disc}L|$.
The polynomial theorem gives
$\alpha_y\Delta^\mu\leq\mathcal B(H)$.
We first bound the torus-level index
$J=[K_T^m:K_T]=\prod_p q_p$.
If $P=\{p:q_p>1\}$, the definition
$\alpha_y=\prod_p\max(1,Bq_p)$ implies
$J\leq B^{-|P|}\alpha_y$.
There are at most $C_0$ primes in $P$ below the cutoff and
at most $\log_2\Delta$ primes ramified in $L$.
At each of the remaining $i$ primes the local index bound is
$q_p\geq c_Up$.  Ordering these primes gives
\[
 \alpha_y\geq (B')^i i!\geq (B'i/e)^i.
\]
If $i>e^2/B'$, the last quantity exceeds $e^i$.
Thus $i\leq e^2/B'+\log\alpha_y$ and
\begin{equation}\label{eq:level-orbit-upper}
 J\leq B^{-(C_0+e^2/B')}
       \alpha_y^{1+\lambda}\Delta^{\lambda/\log2}.
\end{equation}
This is the index estimate immediately following
\cite[Proposition~6.6]{Jaswal}; the argument records all its factors.

Minkowski's ideal-class bound is at most $2^N\sqrt\Delta$.
The number of integral ideals of norm $k$ is at most the
$[L:\Q]$-fold divisor function.  Its elementary convolution
estimate, obtained by summing successively over the first
$[L:\Q]-1$ divisors and using
$\sum_{j\leq X}1/j\leq1+\log X$, gives
\[
 h_L\leq2^N\sqrt\Delta
       (1+N\log2+\tfrac12\log\Delta)^{N-1}.
\]
The torus-class estimate
\cite[Proposition~9.7]{Jaswal} therefore gives
\[
 h(T)\leq N^{r(\log_2N+N)}
       2^{Nr}\Delta^{r/2}
       (1+N\log2+\tfrac12\log\Delta)^R.
\]
Here $h(T)=|T(\Q)\backslash T(\mathbf A_f)/K_T^m|$.

The reflex field $E(T,y)$ is contained in $L$.
Over its compositum with $F$, Shimura reciprocity places the
orbit in $T(\Q)\backslash T(\mathbf A_f)/K_T$;
see \cite[\S2.1]{UllmoYafaev}.
This set has at most $h(T)J$ elements.
Passing from that compositum to $F$ costs a factor at most $N$,
and passing from $y$ to $x$ costs at most $h_f$.
Combining these estimates yields
\[
 |\Gal_Fx|\leq C_{\rm orb}
 \alpha_y^{1+\lambda}\Delta^{r/2+\lambda/\log2}
 (1+N\log2+\tfrac12\log\Delta)^R.
\]
Both $\alpha_y$ and $\Delta$ are at least one.  The definition
of $p_{\rm orb}$ implies that their product on the right is
at most $(\alpha_y\Delta^\mu)^{p_{\rm orb}}$.
Also $\log\Delta\leq\log\mathcal B(H)/\mu$.
These two substitutions give $U(H)$.  Applying the result to
the proper detector image proves the final implication.
\end{proof}

\begin{corollary}\label{cor:centre-forcing}
Assume in addition the Siegel orbit lower bound in
Proposition~\ref{prop:orbit-transfer}, with constant $c_g$.
A relative-Weyl detector of degree at most $D$ at a point $x$
with
\[
 D_x>\left(\frac{[F:\Q]U(H_D)}{c_g}\right)^{4g}
\]
forces $\mathcal D_D=S_0$.
\end{corollary}
\begin{proof}
The lower bound would make the orbit strictly larger than
$U(H_D)$, contradicting Proposition~\ref{prop:orbit-upper}
if the detector image were proper.
\end{proof}

This implication uses a lower and an upper bound for the same
orbit.  It needs no upper bound for an endomorphism-order
discriminant in terms of a torus discriminant alone.


\section{Special cycles on the base and cycles in the fibre}
\label{sec:theta-scope}

This section explains what a theta-lift calculation would have to
compute in order to give the detector used above.
For a unitary domain of signature $(p,q)$, the special cycle
associated with a positive Hermitian subspace of dimension $r$
has codimension $rq$: its domain has dimension $(p-r)q$
instead of $pq$.  Thus a codimension-$n$ positive special cycle
on a signature-$(n,n)$ Shimura variety corresponds to $r=1$.
It is a cycle on the Shimura base, of dimension $n^2-n$.
Its degree depends on its Hermitian datum, level, and the chosen
ample bundle on a specified compactification.

\begin{proposition}[The fibre restriction]
\label{prop:base-cycle-restriction}
Let $\pi:\mathcal A\to S$ be an abelian scheme over a smooth
complex base, let $x\in S(\mathbf C)$, and let
$z\in H^{2n}(S,\mathbf Q)$ with $n>0$.  The restriction of
$\pi^*z$ to $\mathcal A_x$ is zero.  In particular pulling back
a positive-codimension special cycle from the base cannot by
itself give a nonzero Weil class in a fibre.
\end{proposition}

\begin{proof}
The composite $\mathcal A_x\to\mathcal A\to S$ factors through
the point $x$.  Its pullback in degree $2n$ therefore factors
through $H^{2n}(\{x\},\mathbf Q)=0$.
\end{proof}

An incidence correspondence in a Kuga variety could have a
different fibre restriction.  Such a correspondence must be
specified and its nonzero Weil projection proved before its
degree can be identified with a Rallis formula.

The distinction between a value at $3/2$ and a central value is
also a distinction of theta problems.  Haan's formula for the
rank-one to rank-three lift contains
\[
 \frac{L_E(3/2,\operatorname{BC}(\tau)\otimes\gamma^3)}
      {L_F(3,\chi_{E/F})}
 \prod_v Z_v^{\#}(1)
\]
\cite[Theorem~3.4, p.~8]{Haan}.  His equal-rank, rank-one lift
instead contains
\[
 \frac12\frac{L_E(1/2,\operatorname{BC}(\tau)\otimes\gamma)}
                {L_F(1,\chi_{E/F})}
 \prod_v Z_v^{\#}(0)
\]
\cite[Theorem~3.5, p.~9]{Haan}.  The representations, twists,
and local normalizations change along with the ranks.  These
two formulas do not identify the polarization degree of an
unspecified fibre detector.

For comparison, a unitary Hecke character $\chi$ satisfies, at
real $\sigma>1$, the elementary bounds for its Euler product over nonarchimedean factors
\[
 \frac{\zeta_E(2\sigma)}{\zeta_E(\sigma)}
 \leq |L_E(\sigma,\chi)|\leq\zeta_E(\sigma).
\]
At each unramified prime the absolute value of the Euler factor
lies between $(1+N\mathfrak p^{-\sigma})^{-1}$ and
$(1-N\mathfrak p^{-\sigma})^{-1}$; the factor $1$ at a
ramified prime satisfies the same bounds.  Multiplication gives
the displayed inequality.  It applies at $3/2$, without a
Riemann-hypothesis input.  At $1/2$ neither this argument nor
nonvanishing follows.  Control of local factors, integral
normalization, and a correspondence to the fibre would still
be required in either case.

\subsection{Small detectors whose parameter points are not flat}
\label{subsec:doubling-nonflat}

A small cycle on a decomposed fibre need not give the relative
flatness required in Theorem~\ref{thm:flat}.  The following
calculation exhibits that failure on the principal Gaussian
family, including at good split primes.

\begin{proposition}[The coordinate detector on a doubled abelian variety]
\label{prop:doubling-nonflat}
Let $B$ be a principally polarized abelian variety of dimension
$n\geq2$ over an algebraically closed field $k$ of characteristic
zero or characteristic $p\nmid6$.  Put
\[
 A=B_x\times B_y,\qquad I(x,y)=(-y,x),
 \qquad Z=\{0\}\times B_y,
\]
and give $A$ the product principal polarization.
Choose compatible level-three data and let $s$ be the
corresponding point of the smooth Gaussian unitary moduli
scheme of signature $(n,n)$ and hyperbolic type $U^n$.
The cycle $Z$ has degree $n!$ and nonzero Weil projection.
Let $\mathscr H\to S$ be the relative Hilbert scheme for the
Hilbert polynomial of $Z$, using a fixed relatively very ample
power of the polarization, and let $h=[Z]$.
Then
\[
 \dim T_h\mathscr H=n+\frac{n(n+1)}2,
 \qquad
 \mathscr H_s\text{ is smooth of dimension }n\text{ at }h,
 \qquad \dim S=n^2.
\]
The image of $T_h\mathscr H\to T_sS$ has codimension
$n(n-1)/2$.  In particular $\mathscr H\to S$ is not flat at
$h$.  The same holds for the detecting union of Hilbert
components containing $h$.
\end{proposition}

\begin{proof}
We first check the lattice type in characteristic zero.
Choose a symplectic basis $e_j,f_j$ for the principal lattice
of $B$, with $E_B(e_j,f_l)=\delta_{jl}$.  On the product use
\[
 h(v,w)=E_A(Iv,w)+iE_A(v,w).
\]
The vectors $a_j=(e_j,0)$ and $b_j=(0,f_j)$ are a basis over
$\Z[I]$: together with their $I$-multiples they are a
$\Z$-basis of the product lattice.  Their pairings are
$h(a_j,a_l)=h(b_j,b_l)=0$ and
$h(a_j,b_l)=h(b_l,a_j)=\delta_{jl}$.
Thus the Hermitian lattice is $U^n$.  The $i$ and $-i$
eigenspaces of $I$ on the tangent space each have dimension
$n$, so the signature is $(n,n)$.  In positive characteristic
one may lift the principally polarized $B$ and its prime-to-$p$
level data locally to characteristic zero; smoothness of the
Siegel moduli scheme gives this lift.  The doubled construction
therefore has the same prescribed integral type.

Restriction of the polarization to $Z$ is the principal
polarization on $B$, so $\deg Z=\theta_B^n=n!$ by
Riemann--Roch.  For the Weil projection write
$H^1(A)=V_x\oplus V_y$.  After adjoining $i$, projection from
$V_x$ to either $I$-eigenspace is an isomorphism onto that
eigenspace.  The class of $Z$ is a nonzero top exterior form
on $V_x$, and both of its determinant projections are therefore
nonzero.  This argument is valid in Betti cohomology and in
$\ell$-adic cohomology for any auxiliary $\ell\ne p$.
It proves the detector assertion without an ordinary-reduction
hypothesis.

The normal bundle of $Z$ in $A$ is
$\mathcal O_B\otimes\operatorname{Lie}(B_x)$.  Consequently
$h^0(N_{Z/A})=n$ and $h^1(N_{Z/A})=n^2$.
The translation orbit $A/Z$ is an $n$-dimensional closed family
in the Hilbert fibre, and its tangent space is exactly
$H^0(N_{Z/A})$.  Since this is also the Hilbert tangent space,
the local dimension and embedding dimension of the fibre at
$h$ are both $n$.  The fibre is regular, hence smooth over the
algebraically closed field at $h$.

For completeness, the relative tangent sequence is
\begin{equation}
 \label{eq:doubling-hilbert-tangent}
 0\longrightarrow H^0(N_{Z/A})
 \longrightarrow T_h\mathscr H
 \longrightarrow T_sS
 \xrightarrow{\mathrm{ob}} H^1(N_{Z/A}).
\end{equation}
It can be seen directly on an affine cover of $Z$.
Local lifts of its equations in a first-order deformation of
$A$ exist because the embedding is regular.  Their differences
on overlaps form the normal projection of the ambient
Kodaira--Spencer cocycle.  Its class is the obstruction
$\mathrm{ob}$; it vanishes exactly when the local lifts glue,
and the choices of a gluing differ by a global normal section.
This proves exactness and identifies the last map.  This is
also the embedded form of the deformation sequence in
\cite[Lemma~93.10.3, Tag~0E3S]{StacksDeformation}.

Use the principal polarization to identify the tangent space
to the Siegel moduli scheme with symmetric $2n$-by-$2n$
matrices.  This is the Kodaira--Spencer description, valid in
positive characteristic as well; see
\cite[Lemma~12(1), pp.~165--166; \S IV.6,
pp.~178--181]{KoblitzDeformation}.
In the two factor blocks the Gaussian action has matrix
\[
 J=\begin{pmatrix}0&-1_n\\1_n&0\end{pmatrix}.
\]
The unitary tangent vectors are the symmetric matrices fixed
by $M\mapsto JMJ^{\mathsf t}$.  Explicitly they are
\[
 M=\begin{pmatrix}U&V\\-V&U\end{pmatrix},
 \qquad U^{\mathsf t}=U,\qquad V^{\mathsf t}=-V.
\]
There are $n(n+1)/2+n(n-1)/2=n^2$ parameters.
Equivalently, the forgetful map from the smooth unitary moduli
scheme has injective differential, since endomorphisms of an
abelian scheme are rigid; its image lies in this fixed space
and has the same dimension.  Thus the description is exact.
The restriction and normal projection in
\eqref{eq:doubling-hilbert-tangent} send $M$ to its off-diagonal
block $V$, up to a harmless transpose sign.  The obstruction
map has rank $n(n-1)/2$, its kernel has dimension $n(n+1)/2$,
and exactness gives the asserted Hilbert tangent dimension.
All these computations are over $k$; the skew-symmetric count
uses $2\ne0$.

If $\mathscr H\to S$ were flat at $h$, its finite presentation
and geometrically regular fibre there would make it smooth at
$h$.  Its tangent map would be surjective, contradicting the
positive rank $n(n-1)/2$ just computed.  This step uses the
smooth fibre, not tangent failure alone.  Nonzero Weil
projection is locally constant on the chosen cycle-parameter
components, so selecting the detecting union does not alter
the germ at $h$.  The nonflatness persists there.
\end{proof}

\begin{example}[A degree-$24$ detector at the split prime $5$]
\label{ex:doubling-prime-five}
Let $T:y^2=x^3-x$, take $B=T^4$ and $A=B\times B$ with the
Gaussian action of Proposition~\ref{prop:doubling-nonflat}.
The elliptic curve has eight points over $\mathbf F_5$:
the numbers above $x=0,1,2,3,4$ are $1,1,2,2,1$, and there
is one point at infinity.  Its Frobenius trace is $-2$, so it
is ordinary.  The rational prime $5$ splits in $\Q(i)$, the
splitting field of this decomposed CM seed.  After choosing
level-three data over a finite residue-field extension, the
coordinate fourfold is a detector of degree $24$.
Its Hilbert tangent dimension is $14$, its Hilbert fibre
dimension is $4$, and the unitary base dimension is $16$.
The parameter morphism is not flat at this detector.  Thus this
explicit good ordinary split reduction does not satisfy the
relative flatness condition in Theorem~\ref{thm:flat}.
The split seed is not a full relative-Weyl CM point either:
its endomorphism algebra contains the matrix algebra coming
from the repeated factor.
\end{example}


\section{What remains to construct}

The unresolved geometric step for the principal Gaussian
eightfold family is a detector of degree $D$ together with
either the discriminant inequality \eqref{eq:forcing} (or the
smaller threshold \eqref{eq:hybrid-detector-gate} under its height
hypotheses), or flatness of its parameter point over the full moduli
base.
The exact semiregularity criterion in \cite{Bounded} is a
possible sufficient condition: a coherently equivariant gluable
perfect object of pure Weil character would have to have
invariant second self-Ext dimension $56$. The numerical
criterion alone constructs neither the object nor a cycle.

For the principal level-three families,
Theorem~\ref{thm:explicit-theta-chow-bound} supplies a completely
specified integer $H_g(D)$, including the comparison with the
Baily--Borel polarization.  The minimum of the two thresholds in
Corollary~\ref{cor:hybrid-detector} can therefore be evaluated from
$D$ and the explicitly stated analytic inputs.
The arithmetic certificate of
Theorem~\ref{thm:arithmetic} is an algorithm from equations,
not an evaluated integer for the entire Gaussian eightfold
family. Pointwise CM algebraicity and pointwise Tate algebraicity
do not remove these requirements.

\section*{Method, checks and records}
The mathematics and text were produced by GPT-6 Astra (ChatGPT
Pro mode). Same-model independent re-proofs are recorded separately
from proof-aware checks and finite computations. The earlier
materials record checks by Claude Opus 5.5 at high and maximum
effort; their scope is not enlarged here. The current interface
does not expose a reliable reasoning-effort setting. No human
has checked this work. The original text and programs are
dedicated to the public domain under CC0 1.0.

The \emph{Algebraic cycles and motives workbench},
\url{\workbenchurl}, gives the exact source locators, programs,
certificates and check records. The permanent entry point for
the records is \url{\recordurl}.

\begin{thebibliography}{99}
\bibitem{Bounded} \emph{The Hodge conjecture for abelian varieties as
a uniform bound over finite fields}. Companion article.
Family criterion; sections on cycle parameter spaces, lifting,
Weil extraction and one finite-field detector. \url{\workbenchurl}.
\bibitem{BrunjesSerpe} L.~Br\"unjes and C.~Serp\'e,
\emph{Uniform bounds and ultraproducts of cycles},
arXiv:0901.4853, Theorem~3.4, p.~8.
\url{https://arxiv.org/abs/0901.4853}.
\bibitem{PST} J.~Pila, A.~Shankar and J.~Tsimerman, with an
appendix by H.~Esnault and M.~Groechenig,
\emph{Canonical heights on Shimura varieties and the
Andr\'e--Oort conjecture}, arXiv:2109.08788, Theorem~1.1.
\url{https://arxiv.org/abs/2109.08788}.
\bibitem{Jaswal} M.~Jaswal, \emph{An effective Andr\'e--Oort
theorem for Shimura varieties}, Hexagon 2611.00339v3,
12 November 2026. Theorem~1.1(ii), Corollary~1.2,
Remark~8.2. Unrefereed.
\url{https://hexagonmath.org/2611.00339v3}.
\bibitem{CuretonThorner} D.~Cureton and J.~Thorner,
\emph{A quasi-Riemann hypothesis for Dedekind zeta functions},
Hexagon 2610.00200v2, Theorem~1.1. Unrefereed.
\url{https://hexagonmath.org/2610.00200v2}.
\bibitem{Hu} B.~Hu, \emph{A seven-eighths zero-free half-plane
for unitary Weil $L$-functions}, Hexagon 2610.00216v2,
Theorem~1.2 and Corollary~9.6. Unrefereed.
\url{https://hexagonmath.org/2610.00216v2}.
\bibitem{Alpoge} L.~Alp\"oge, \emph{Announcement on the Riemann
hypothesis}, 20 November 2026, p.~1 and p.~4.
\url{https://alpo.ge/remark.pdf}.
\bibitem{MilneReport} J.~S.~Milne, \emph{On the Hodge, Tate,
and standard conjectures for abelian varieties},
Hexagon 2611.00244v1, 9 November 2026, Theorems~3.1 and~3.5.
Report dependent on the OpenAI preprints of 6 October 2026.
\url{https://hexagonmath.org/2611.00244v1}.
\bibitem{ChaiOort} C.-L.~Chai and F.~Oort,
\emph{Abelian varieties isogenous to a Jacobian},
Ann. of Math. 176 (2012), 589--635.
Definitions~5.3--5.4, Remark~5.8, Lemma~6.3.
\url{https://doi.org/10.4007/annals.2012.176.1.11}.
\bibitem{SGA45} P.~Deligne, \emph{Cohomologie \'etale},
SGA~$4\frac12$, LNM~569 (1977), Cycle, \S2.3.1--2.3.8;
erratum to SGA~4 XVI, Corollary~2.2.
\bibitem{BachSorenson} E.~Bach and J.~Sorenson,
\emph{Explicit bounds for primes in residue classes},
Math. Comp. 65 (1996), 1717--1735, Theorem~5.3,
pp.~1731--1732.
\url{https://doi.org/10.1090/S0025-5718-96-00763-6}.

\bibitem{UllmoYafaev} E.~Ullmo and A.~Yafaev,
\emph{Galois orbits and equidistribution of special subvarieties:
towards the Andr\'e--Oort conjecture}, Ann. of Math. 180 (2014),
823--865, \S2.1 and Lemmas~2.4--2.8, Proposition~2.9.
\url{https://www.ihes.fr/~/ullmo/publications/UllmoYafaevSept2013.pdf}.
\bibitem{ConradST} B.~Conrad,
\emph{The Shimura--Taniyama formula}, seminar notes (2004),
Theorem~2.1, p.~3.
\url{https://math.stanford.edu/~conrad/vigregroup/vigre04/stformula.pdf}.
\bibitem{Haan} J.~Haan, \emph{The Bessel period of $U(3)$ and $U(2)$
involving a non-tempered representation}, arXiv:1403.5061v7,
Sections~3.3--3.4, Theorems~3.4--3.5, pp.~8--9.
\url{https://arxiv.org/abs/1403.5061}.

\bibitem{EdixhovenYafaev}
B.~Edixhoven and A.~Yafaev, \emph{Subvarieties of Shimura varieties},
Ann. of Math. \textbf{157} (2003), 621--645.
Used author version May 29, 2001: Proposition~7.4.3, p.~23.
\url{https://websites.math.leidenuniv.nl/edixhoven/public_html_rennes/publications/anoort.pdf}.

\bibitem{MilneShimura}
J.~S.~Milne, \emph{Introduction to Shimura Varieties},
revised September 16, 2017.
Used: Theorem~5.17, p.~59; connected-component reciprocity,
\S13, formula~(64), p.~119; the Siegel cyclotomic action, p.~125.
\url{https://www.jmilne.org/math/xnotes/svi.pdf}.


\bibitem{MilneAV}
J.~S.~Milne, \emph{Abelian Varieties},
Chapter V of \emph{Arithmetic Geometry} (Storrs, 1984),
103--150, Springer, 1986; corrected author text, January 2, 2022.
Used: Remark~7.3, p.~12, recording the third-power very-ampleness theorem.
\url{https://www.jmilne.org/math/xnotes/AVs.pdf}.
\bibitem{DeligneAH82WO}
P.~Deligne (notes by J.~S.~Milne),
\emph{Hodge cycles on abelian varieties},
in P.~Deligne, J.~S.~Milne, A.~Ogus and K.-y.~Shih,
\emph{Hodge Cycles, Motives, and Shimura Varieties},
Lecture Notes in Mathematics 900, Springer, 1982, pp.~9--100.
Locators in the TeXed version of 1 October 2018:
Proposition~2.9(b), pp.~18--19, and Main Theorem~2.11, p.~19.
\url{https://jmilne.org/math/Documents/Deligne82.pdf}.

\bibitem{Tate66WO}
J.~Tate,
\emph{Endomorphisms of abelian varieties over finite fields},
Inventiones Mathematicae \textbf{2} (1966), 134--144.
Main Theorem, p.~134, and Lemma~1, p.~135.
\url{https://pazuki.perso.math.cnrs.fr/index_fichiers/Tate66.pdf}.

\bibitem{JonesRoberts14WO}
J.~W.~Jones and D.~P.~Roberts,
\emph{The tame-wild principle for discriminant relations for number fields},
Algebra \& Number Theory \textbf{8} (2014), no.~3, 609--645.
DOI: \texttt{10.2140/ant.2014.8.609}.
Section~3C, equations (3-5)--(3-7), pp.~614--615;
Theorem~7.3, p.~634.
\url{https://msp.org/ant/2014/8-3/ant-v8-n3-p04-p.pdf}.

\bibitem{Wei} X.~Wei, \emph{Smallness of Faltings heights of CM
abelian varieties}, J. Number Theory 242 (2023), 63--94.
Author version arXiv:2111.00617v1, Lemma~4.2, pp.~12--14;
Proposition~5.3, pp.~19--20; height normalization, p.~8.
\url{https://arxiv.org/abs/2111.00617}.
\bibitem{GR14} E.~Gaudron and G.~R\'emond,
\emph{Th\'eor\`eme des p\'eriodes et degr\'es minimaux d'isog\'enies},
Comment. Math. Helv. 89 (2014), 343--403.
Section~2.3, pp.~351--352; Corollary~8.4, p.~400.
\url{https://doi.org/10.4171/CMH/322}.
\bibitem{GR23} E.~Gaudron and G.~R\'emond,
\emph{Nouveaux th\'eor\`emes d'isog\'enie},
M\'em. Soc. Math. France 176 (2023).
Notations~1.1--1.2; Theorems~1.5(2), 1.8; Proposition~13.7.
Author version September 2022, pp.~11, 13, 15, 89--90.
\url{https://lmbp.uca.fr/~gaudron/art20.pdf}.
\bibitem{YuanZhang} X.~Yuan and S.-W.~Zhang,
\emph{On the averaged Colmez conjecture},
Ann. of Math. 187 (2018), 533--638.
Theorem~1.1, p.~535; normalization, p.~534.
\url{https://doi.org/10.4007/annals.2018.187.2.4}.
\bibitem{AGHMP} F.~Andreatta, E.~Z.~Goren, B.~Howard and
K.~Madapusi Pera, \emph{Faltings heights of abelian varieties
with complex multiplication}, Ann. of Math. 187 (2018), 391--531.
Theorem~A, p.~392; normalization, Section~9.1.
\url{https://doi.org/10.4007/annals.2018.187.2.3}.
\bibitem{Tsimerman} J.~Tsimerman, \emph{The Andr\'e--Oort conjecture
for $\mathcal A_g$}, Ann. of Math. 187 (2018), 379--390.
Theorem~1.3, p.~380; Theorems~4.1 and~5.1.
\url{https://doi.org/10.4007/annals.2018.187.2.2}.

\bibitem{Mumford1966I}
D.~Mumford,
\emph{On the equations defining abelian varieties. I},
Invent. Math. \textbf{1} (1966), 287--354.
\url{https://pazuki.perso.math.cnrs.fr/index_fichiers/Mumford66.pdf}.

\bibitem{Mumford1967II}
D.~Mumford,
\emph{On the equations defining abelian varieties. II},
Invent. Math. \textbf{3} (1967), 75--135.
\url{https://doi.org/10.1007/BF01389741}.
Author copy:
\url{https://www.dam.brown.edu/people/mumford/alg_geom/papers/1967a--EqDefAVIIcomp-GT.pdf}.

\bibitem{Mumford1967III}
D.~Mumford,
\emph{On the equations defining abelian varieties. III},
Invent. Math. \textbf{3} (1967), 215--244.
Author copy:
\url{https://www.dam.brown.edu/people/mumford/alg_geom/papers/1967b--EqDefAVIIIcomp-GT.pdf}.

\bibitem{FreitagSalvatiManni2016}
E.~Freitag and R.~Salvati Manni,
\emph{On the variety associated to the ring of theta constants in genus 3},
arXiv:1606.04468v1 (2016).
\url{https://arxiv.org/abs/1606.04468}.
The references here to the general theta map are to pp.~1--3,
not to a claim of normality in genus four or eight.

\bibitem{GuerraPirola2009}
L.~Guerra and G.~P.~Pirola,
\emph{On the finiteness theorem for rational maps on a variety of general type},
Collect. Math. \textbf{60} (2009), 261--276.
\url{https://doi.org/10.1007/BF03191371}.
Page locators in the text refer to arXiv:0807.2980v2,
\url{https://arxiv.org/abs/0807.2980}.

\bibitem{LeFourn2019}
S.~Le Fourn,
\emph{A tubular variant of Runge's method in all dimensions,
with applications to integral points on Siegel modular varieties},
Algebra Number Theory \textbf{13} (2019), 159--209.
\url{https://doi.org/10.2140/ant.2019.13.159}.
The marking discussion used here is in \S8A, pp.~198--200.

\bibitem{DawJavanpeykarKuhne2020}
C.~Daw, A.~Javanpeykar, and L.~K\"uhne,
\emph{Effective estimates for the degrees of maximal special subvarieties},
Selecta Math. (N.S.) \textbf{26} (2020), article~2.
\url{https://doi.org/10.1007/s00029-019-0528-1}.

\bibitem{StacksDeformation} The Stacks project,
Lemma~93.10.3, Tag~0E3S.
\url{https://stacks.math.columbia.edu/tag/0E3S}.
\bibitem{KoblitzDeformation} N.~Koblitz,
\emph{$p$-adic variation of the zeta-function over families of
varieties defined over finite fields}, Compositio Math.
\textbf{31} (1975), 119--218, Lemma~12(1), pp.~165--166;
\S IV.6, pp.~178--181.
\url{https://www.numdam.org/article/CM_1975__31_2_119_0.pdf}.


\bibitem{FGA} A.~Grothendieck,
\emph{Techniques de construction et th\'eor\`emes d'existence
en g\'eom\'etrie alg\'ebrique. IV. Les sch\'emas de Hilbert},
S\'eminaire Bourbaki, Expos\'e~221 (1961),
Theorem~3.2, p.~260.
\url{https://numdam.org/item/SB_1960-1961__6__249_0/}.
\bibitem{SGA6} P.~Berthelot, A.~Grothendieck and L.~Illusie, eds.,
\emph{Th\'eorie des intersections et th\'eor\`eme de Riemann--Roch},
Lecture Notes in Math. \textbf{225}, Springer, 1971,
Expos\'e~XIII (S.~Kleiman), Definition~6.9 and Corollary~6.11,
pp.~658--660.

\end{thebibliography}
\end{document}
